\begingroup
\setlength{\emergencystretch}{3em}
\newcommand{\chapterSevenDash}{{\fontfamily{lmr}\selectfont\textemdash}}
\chapter{Đóng gói, triển khai và vận hành}
\label{chap:dong-goi-van-hanh}

Chương này trình bày cách đóng gói, cấu hình, triển khai và vận hành hệ thống RAG\_TS. Chương này dựa trên mã nguồn của hai Git submodule \texttt{backend\_\allowbreak{}rag} (NestJS 11, TypeScript 5.7) và \texttt{frontend\_\allowbreak{}rag} (React 19, Vite 6, React Three Fiber), tại nhánh \texttt{Minh02112026} của kho \texttt{RAG\_\allowbreak{}TS}.

RAG\_TS được triển khai on-premises. Thay vì gọi API của nhà cung cấp LLM, nhóm vận hành mô hình embedding, chỉ mục vector và mô hình sinh văn bản trên hạ tầng nội bộ. Vì vậy, chi phí suy luận phụ thuộc vào tài nguyên máy chủ. Một số lỗi cấu hình cũng khó phát hiện qua mã trạng thái HTTP: hệ thống có thể trả lời bình thường nhưng truy hồi sai tài liệu. Các mục sau trình bày cách kiểm tra cấu hình khi khởi động, đóng gói ứng dụng bằng multi-stage build và kiểm tra luồng hỏi đáp sau khi triển khai.


\section{Chuẩn bị môi trường triển khai}
\label{sec:chuan-bi-moi-truong}

\subsection{Quy mô tài nguyên và cách xác định yêu cầu phần cứng}

Một lượt hỏi đáp trong RAG\_TS đòi hỏi năm thành phần phải cùng sẵn sàng: chỉ mục vector Qdrant phải truy vấn được, mô hình embedding phải nạp được để biến câu hỏi thành vector cùng không gian với các đoạn văn bản, mô hình ngôn ngữ phải nạp được để sinh câu trả lời, cơ sở dữ liệu MongoDB phải lưu được tài khoản và lịch sử hội thoại, và giao diện web phải được phục vụ tĩnh. Việc trình duyệt tải được giao diện chỉ xác nhận một phần nhỏ của hệ thống đã khởi động, vì giao diện được build thành các tệp tĩnh và không phụ thuộc vào API.

Điểm cần lưu ý khi chọn phần cứng là kích thước tệp trọng số trên đĩa không phản ánh lượng bộ nhớ cần thiết khi xử lý một lượt hỏi đáp. Bản lượng tử hóa Q4\_K\_M của Qwen2.5-7B-Instruct có dung lượng khoảng 4,68 GB, nhưng khi mô hình được nạp vào bộ nhớ và phải lưu thêm ngữ cảnh sinh văn bản, tổng bộ nhớ cần thiết lớn hơn đáng kể. Tài liệu vận hành của dự án ghi nhận mức khoảng 6 GB RAM cho container chạy mô hình 7B trên CPU. Khi chọn thiết bị, cách xác định đáng tin cậy nhất là đo trực tiếp trên chính máy đích: đo bộ nhớ cư trú và thời gian từ khi gửi yêu cầu đến khi nhận token đầu tiên (time to first token, TTFT) với một câu hỏi đại diện cho tải thực tế.

Tương ứng với quy mô dữ liệu của hệ thống, mô hình embedding Nomic Embed Text v1.5 có dung lượng khoảng 550 MB và được tải tự động từ Hugging Face ở lần chạy đầu tiên, sau đó được tái sử dụng từ bộ nhớ đệm.

\subsection{Thành phần dữ liệu đầu vào và điều kiện phiên bản}

Hệ thống phân biệt rõ hai loại tài nguyên: dữ liệu do dự án sở hữu và dữ liệu do pipeline bên ngoài xuất ra.

Thư mục \texttt{model/\allowbreak{}} chứa trọng số mô hình ngôn ngữ ở định dạng GGUF, được gắn vào container backend dưới đường dẫn \texttt{/\allowbreak{}models} ở chế độ chỉ đọc. Thư mục \texttt{data/\allowbreak{}} chứa các tệp mà pipeline khai phá xuất ra, được gắn vào \texttt{/\allowbreak{}data} cũng ở chế độ chỉ đọc. Cụ thể, \texttt{WORLD\_\allowbreak{}DATA\_\allowbreak{}DIR} trỏ tới \texttt{/\allowbreak{}data/\allowbreak{}world} chứa bốn tệp phục vụ cây tri thức 3D \chapterSevenDash{} \texttt{tree.\allowbreak{}json}, \texttt{centroids\_\allowbreak{}l1.\allowbreak{}json}, \texttt{centroids\_\allowbreak{}l2.\allowbreak{}json} và \texttt{categories.\allowbreak{}json} \chapterSevenDash{} còn \texttt{MINING\_\allowbreak{}DIR} trỏ tới \texttt{/\allowbreak{}data/\allowbreak{}gold/\allowbreak{}mining} chứa \texttt{graph\_\allowbreak{}coauthorship.\allowbreak{}json} (điểm PageRank đồng tác giả) và \texttt{association\_\allowbreak{}rules.\allowbreak{}json} (luật kết hợp FP-Growth).

Điểm cần lưu ý về các tệp dữ liệu tùy chọn: nếu thư mục world hoặc mining vắng mặt, tiến trình backend vẫn khởi động bình thường. Một số chức năng sẽ bị hạn chế \chapterSevenDash{} sự kiện \texttt{route} trong ngữ cảnh truy hồi sẽ có \texttt{steps:\allowbreak{} []}, tính năng xếp hạng lại theo PageRank không hoạt động, và trường \texttt{category\_\allowbreak{}index} trả về giá trị \texttt{-\allowbreak{}1}. Đây là một quyết định thiết kế có chủ ý: hệ thống vẫn cung cấp được chức năng hỏi đáp cơ bản ngay cả khi pipeline chưa xuất xong dữ liệu, thay vì từ chối khởi động. Người vận hành cần biết điều này để không nhầm một thiếu sót dữ liệu với một lỗi triển khai.

Điều kiện bắt buộc về phiên bản: mã bài báo trong chỉ mục vector và trong dữ liệu khai phá phải cùng thuộc một phiên bản. Nếu hai nguồn lệch phiên bản, kết quả truy hồi sẽ trả về các đoạn văn bản không khớp với cây tri thức đang hiển thị trên giao diện, dẫn tới tình huống trích dẫn mở được bài báo sai.

\subsection{Yêu cầu về phần mềm nền tảng}

Backend yêu cầu Node.js 22 trở lên (do \texttt{node-\allowbreak{}llama-\allowbreak{}cpp} và các thư viện native cần ABI tương thích), frontend yêu cầu Node.js 20 trở lên. Hệ thống container hoá cần Docker kèm Docker Compose v2.

Đối với môi trường phát triển, hệ thống MongoDB từ xa đặt sau Cloudflare Access nên không thể truy cập trực tiếp bằng giao thức TCP thông thường. Người vận hành phải chạy một tiến trình \texttt{cloudflared access tcp} như một daemon cục bộ:

\begin{minted}[fontsize=\footnotesize,breaklines,breakanywhere]{bash}
cloudflared access tcp --hostname mongodb.minio-quangmanh.io.vn --url localhost:27018
\end{minted}

Cổng 27018 được chọn có chủ ý để tránh xung đột với MongoDB cục bộ ở cổng 27017. Chỉ mục vector Qdrant được phục vụ qua HTTPS nên không cần đường hầm tương tự.

Điểm vận hành đáng lưu ý nhất của kiến trúc này là: mọi lượt khứ hồi tới MongoDB từ xa đều đi qua đường hầm và tiêu tốn khoảng 15–25 ms. Người vận hành cần duy trì tiến trình \texttt{cloudflared} trong suốt thời gian ứng dụng chạy; khi tiến trình này tắt, toàn bộ chức năng cần ghi vào cơ sở dữ liệu sẽ thất bại dù các thành phần khác vẫn hoạt động.


\section{Đóng gói ứng dụng và các thành phần hệ thống}
\label{sec:dong-goi-ung-dung}

\subsection{Dockerfile của backend: hai stage build và runtime}

\texttt{backend\_\allowbreak{}rag/\allowbreak{}Dockerfile} dùng hai giai đoạn trên nền \texttt{node:\allowbreak{}22-\allowbreak{}bookworm-\allowbreak{}slim}.

Build stage (\texttt{build}) thực hiện bốn bước theo thứ tự để tối ưu bộ nhớ đệm của Docker: sao chép \texttt{package.\allowbreak{}json} và \texttt{package-\allowbreak{}lock.\allowbreak{}json} trước, chạy \texttt{npm ci} để cài đặt theo tệp khóa, sao chép riêng các tệp cấu hình biên dịch (\texttt{tsconfig.\allowbreak{}json}, \texttt{tsconfig.\allowbreak{}build.\allowbreak{}json}, \texttt{nest-\allowbreak{}cli.\allowbreak{}json}), rồi mới sao chép mã nguồn \texttt{src}. Nhờ thứ tự này, Docker tái sử dụng layer \texttt{npm ci} khi chỉ mã nguồn thay đổi và các tệp khai báo phụ thuộc vẫn giữ nguyên. Sau đó \texttt{npm run build} biên dịch TypeScript sang thư mục \texttt{dist} và \texttt{npm prune -\allowbreak{}-\allowbreak{}omit=\allowbreak{}dev} loại bỏ các gói chỉ phục vụ phát triển.

Runtime stage chỉ chứa các thành phần cần để chạy backend: \texttt{node\_\allowbreak{}modules} đã được cắt tỉa (chuyển quyền sở hữu cho tài khoản \texttt{node} ngay trong lệnh \texttt{COPY -\allowbreak{}-\allowbreak{}chown}), thư mục \texttt{dist} đã biên dịch, và \texttt{package.\allowbreak{}json}. Runtime image đặt \texttt{NODE\_\allowbreak{}ENV=\allowbreak{}production} và \texttt{PORT=\allowbreak{}3000}, chạy dưới tài khoản không phải root.

Hai chi tiết nhỏ nhưng quan trọng về mặt vận hành:

Thứ nhất, runtime image tạo sẵn thư mục bộ nhớ đệm Hugging Face tại \texttt{node\_\allowbreak{}modules/\allowbreak{}@\allowbreak{}huggingface/\allowbreak{}transformers/\allowbreak{}.\allowbreak{}cache} và chuyển quyền sở hữu cho tài khoản \texttt{node}. Điều này là cần thiết vì thư mục này được gắn thêm dưới dạng volume \texttt{hf-\allowbreak{}cache} trong Compose; nếu Docker tạo volume trên một thư mục chưa tồn tại trong image, volume sẽ thuộc sở hữu root và tiến trình chạy dưới tài khoản \texttt{node} sẽ không ghi được.

Thứ hai, Dockerfile khai báo sẵn lệnh \texttt{HEALTHCHECK} gọi \texttt{http:\allowbreak{}/\allowbreak{}/\allowbreak{}127.\allowbreak{}0.\allowbreak{}0.\allowbreak{}1:\allowbreak{}3000/\allowbreak{}api/\allowbreak{}v1/\allowbreak{}health/\allowbreak{}live} với chu kỳ 15 giây, thời gian chờ 3 giây và thời gian khởi động 30 giây. Endpoint \texttt{live} trả HTTP 200 khi tiến trình còn trả lời mà không gọi bất kỳ phụ thuộc nào, nên nó phù hợp cho mục đích kiểm tra trạng thái hoạt động của container.

\subsection{Dockerfile của frontend: build frontend và phục vụ tệp tĩnh bằng Nginx}

\texttt{frontend\_\allowbreak{}rag/\allowbreak{}Dockerfile} có hai giai đoạn với vai trò khác hẳn backend.

Giai đoạn \texttt{builder} dùng \texttt{node:\allowbreak{}20}, chạy \texttt{npm install}, sao chép toàn bộ mã nguồn và \texttt{npm run build}. Lệnh build trong \texttt{package.\allowbreak{}json} là \texttt{tsc -\allowbreak{}b \&\& vite build}, tức là kiểm tra kiểu dữ liệu bằng TypeScript trước rồi mới đóng gói bằng Vite. Giai đoạn \texttt{nginx:\allowbreak{}alpine} chỉ chứa kết quả đóng gói tĩnh trong \texttt{/\allowbreak{}usr/\allowbreak{}share/\allowbreak{}nginx/\allowbreak{}html}.

Một dòng quan trọng cần lưu ý là \texttt{RUN rm /\allowbreak{}etc/\allowbreak{}nginx/\allowbreak{}conf.\allowbreak{}d/\allowbreak{}default.\allowbreak{}conf}. Lệnh này xoá cấu hình máy chủ mặc định của Nginx, vốn lắng nghe ở cổng 80 với \texttt{root} trỏ tới thư mục mặc định. Nếu giữ lại, container sẽ phục vụ sai nội dung và không nhận được khối cấu hình \texttt{server} do người vận hành cung cấp. Vì vậy, bản triển khai bắt buộc phải cung cấp cấu hình Nginx riêng, điều được trình bày ở Mục 7.6.

Tuy nhiên trong bản triển khai thực tế của dự án, frontend \textbf{không được đóng gói thành Docker image chứa Nginx riêng}. Hai tệp Compose ở thư mục gốc kho thay vào đó dùng image \texttt{nginx:\allowbreak{}alpine} thuần và ánh xạ thư mục chứa bản build trên máy host:

\begin{minted}[fontsize=\footnotesize,breaklines,breakanywhere]{yaml}
- ./frontend-rag/dist:/usr/share/nginx/html
\end{minted}

Cách này có hai hệ quả. Thứ nhất, việc triển khai lại giao diện chỉ cần build lại frontend rồi tải lại container Nginx, không cần build lại image \chapterSevenDash{} nhanh hơn đáng kể so với cách đóng gói tệp vào image. Thứ hai, cấu hình Nginx do người vận hành cung cấp thay vì nằm trong image, nên một bản cấu hình duy nhất phục vụ được cho cả môi trường phát triển lẫn production, chỉ khác nhau ở tệp ánh xạ. Đổi lại, thư mục \texttt{dist/\allowbreak{}} phải được build sẵn trên máy chủ trước khi khởi động Compose, và \texttt{dist/\allowbreak{}} cũng nằm trong \texttt{.\allowbreak{}gitignore}.

Dockerfile của frontend mô tả cách build giao diện và phục vụ các tệp tĩnh bằng Nginx. Trong CI, có thể dùng Dockerfile này để tạo một image độc lập. Khi triển khai bằng Compose ở thư mục gốc, Nginx phục vụ cùng bản build từ thư mục được mount trên host.

\subsection{Cấu trúc dịch vụ trong Docker Compose}

Tệp \texttt{backend\_\allowbreak{}rag/\allowbreak{}docker-\allowbreak{}compose.\allowbreak{}yml} định nghĩa ba dịch vụ và ba volume.

Dịch vụ \texttt{mongo} dùng image \texttt{mongo:\allowbreak{}7} với lệnh \texttt{['-\allowbreak{}-\allowbreak{}replSet', 'rs0', '-\allowbreak{}-\allowbreak{}bind\_\allowbreak{}ip\_\allowbreak{}all']}. Việc MongoDB chạy ở chế độ replica set một nút tên \texttt{rs0} là điều kiện tiên quyết cho các giao dịch (transaction) khi ghi lịch sử hội thoại. Cổng ánh xạ ra máy chủ là \texttt{\$\{RAGTS\_\allowbreak{}MONGO\_\allowbreak{}PORT:\allowbreak{}-\allowbreak{}27018\}:\allowbreak{}27017}, cổng 27018 vẫn được giữ để tránh xung đột với MongoDB cục bộ. Khối \texttt{healthcheck} của dịch vụ này là điểm thiết kế đáng chú ý: câu lệnh kiểm tra thử gọi \texttt{rs.\allowbreak{}status().\allowbreak{}ok} và nếu gặp ngoại lệ thì khởi tạo replica set ngay trong lần chạy đầu tiên:

\begin{minted}[fontsize=\footnotesize,breaklines,breakanywhere]{text}
try { rs.status().ok } catch (e) { rs.initiate({ _id: 'rs0', members: [{ _id: 0, host: 'mongo:27017' }] }).ok }
\end{minted}

Cơ chế này giúp container MongoDB khởi động được trên một volume trống mà không cần thao tác thủ công nào. Với chu kỳ 5 giây và tối đa 30 lần thử lại, Compose sẽ chờ tối đa 150 giây để MongoDB đạt trạng thái healthy.

Dịch vụ \texttt{qdrant} dùng image \texttt{qdrant/\allowbreak{}qdrant:\allowbreak{}v1.\allowbreak{}15.\allowbreak{}4} nhưng thuộc Compose profile \texttt{standalone}. Cơ chế profile này cho phép chạy một bản Qdrant rỗng chỉ khi được yêu cầu tường minh bằng \texttt{docker compose -\allowbreak{}-\allowbreak{}profile standalone up}, trong khi lần chạy thông thường chỉ khởi động MongoDB và ứng dụng. Đây là sự phân tách có chủ đích: theo quy ước bất biến của hệ thống, chỉ mục vector \texttt{scientific\_\allowbreak{}papers\_\allowbreak{}gold} do pipeline bên ngoài ghi vào, và backend chỉ đọc. Bản Qdrant rỗng trong profile \texttt{standalone} chỉ dành cho việc chẩn đoán và chạy thử độc lập, kèm tệp \texttt{scripts/\allowbreak{}seed-\allowbreak{}qdrant.\allowbreak{}ts} nạp chỉ mục giả.

Dịch vụ \texttt{app} build từ thư mục hiện tại và ánh xạ cổng \texttt{\$\{RAGTS\_\allowbreak{}PORT:\allowbreak{}-\allowbreak{}3000\}:\allowbreak{}3000}. Khối \texttt{depends\_\allowbreak{}on} với điều kiện \texttt{condition:\allowbreak{} service\_\allowbreak{}healthy} bảo đảm ứng dụng chỉ khởi động sau khi MongoDB đã vượt qua health check, tránh tình trạng ứng dụng khởi động trước và liên tục thử kết nối lại.

Ba volume được định nghĩa là \texttt{mongo-\allowbreak{}data} (lưu trữ cố định của MongoDB), \texttt{qdrant-\allowbreak{}data} và \texttt{hf-\allowbreak{}cache} (bộ nhớ đệm mô hình embedding).

\subsection{Kiến trúc hai tầng của Compose}

Một điểm cần làm rõ về cấu trúc triển khai: hệ thống có \textbf{hai tầng Compose độc lập} chứ không phải một.

Tầng thứ nhất là \texttt{backend\_\allowbreak{}rag/\allowbreak{}docker-\allowbreak{}compose.\allowbreak{}yml}, chịu trách nhiệm về hạ tầng dữ liệu và dịch vụ API: MongoDB với replica set, Qdrant rỗng theo profile, và container ứng dụng. Tầng này không có Nginx.

Tầng thứ hai là cặp \texttt{docker-\allowbreak{}compose.\allowbreak{}yaml} và \texttt{docker-\allowbreak{}compose.\allowbreak{}prod.\allowbreak{}yaml} ở thư mục gốc kho, chịu trách nhiệm về tầng tiếp nhận lưu lượng: Nginx và một instance backend, với \texttt{nginx.\allowbreak{}dev.\allowbreak{}conf} hoặc \texttt{nginx.\allowbreak{}prod.\allowbreak{}conf} tương ứng. Tầng này không có MongoDB.

Việc tách cấu hình Nginx khỏi tầng dữ liệu cho phép quản lý cấu hình Nginx và chứng thư TLS độc lập với vòng đời của cơ sở dữ liệu, và bản triển khai production có thể dùng một backend đã có sẵn dữ liệu thay vì dựng lại toàn bộ stack. Đổi lại, người vận hành phải nhớ khởi động hai tệp Compose theo đúng thứ tự, và tên service \texttt{backend-\allowbreak{}rag} phải khớp giữa tầng biên và tầng dữ liệu.

Cần lưu ý rằng cấu hình ở tầng biên sử dụng cổng và tên service riêng (\texttt{backend-\allowbreak{}rag}, cổng 8000), khác với tầng dữ liệu (\texttt{app}, cổng 3000, \texttt{mongo}, cổng 27018). Sự khác biệt này phản ánh một giai đoạn phát triển khác nhau trong lịch sử dự án: tầng biên được thiết kế khi backend còn dùng cổng 8000 và LanceDB, còn tầng dữ liệu hiện tại đã chuyển sang Qdrant và cổng 3000. Khi triển khai production thực tế, hai tầng này cần được đồng bộ lại về cùng một tên service và cổng.

\subsection{Quy ước đặt tên biến môi trường trong Compose}

Compose nhận các biến có tiền tố \texttt{RAGTS\_\allowbreak{}} từ máy chủ rồi ánh xạ sang tên biến mà ứng dụng đọc trong container. Ví dụ \texttt{RAGTS\_\allowbreak{}QDRANT\_\allowbreak{}URL} xác định giá trị của \texttt{QDRANT\_\allowbreak{}URL}, còn \texttt{RAGTS\_\allowbreak{}MODEL\_\allowbreak{}FILE} xác định tên tệp mô hình dưới \texttt{/\allowbreak{}models}. Nguyên nhân của quy ước này là xung đột tên biến: tệp \texttt{.\allowbreak{}env} ở thư mục gốc \texttt{backend\_\allowbreak{}rag} dùng cho chạy trực tiếp ngoài Docker (nơi \texttt{MONGO\_\allowbreak{}URI} trỏ tới \texttt{localhost:\allowbreak{}27018} qua đường hầm Cloudflare), trong khi trong container, \texttt{MONGO\_\allowbreak{}URI} phải trỏ tới tên dịch vụ \texttt{mongo:\allowbreak{}27017}. Nếu dùng chung một tập tên biến, môi trường phát triển và môi trường container sẽ ghi đè lẫn nhau. Tiền tố \texttt{RAGTS\_\allowbreak{}} tách hai không gian này, trong khi Compose tự đọc tệp \texttt{.\allowbreak{}env} để lấy các biến đó.

\subsection{Kết nối dịch vụ bên ngoài và cấu hình bind mount chỉ đọc}

Compose khai báo hai bind mount chỉ đọc cho mô hình và dữ liệu:

\begin{minted}[fontsize=\footnotesize,breaklines,breakanywhere]{yaml}
volumes:
  - ../model:/models:ro
  - ../data:/data:ro
  - hf-cache:/app/node_modules/@huggingface/transformers/.cache
extra_hosts:
  - 'host.docker.internal:host-gateway'
\end{minted}

Cờ \texttt{:\allowbreak{}ro} (chỉ đọc) trên hai thư mục đầu tiên thể hiện nguyên tắc bất biến của hệ thống: dữ liệu đầu vào và trọng số không bao giờ bị sửa bởi tiến trình ứng dụng. Điều này không chỉ là biện pháp an toàn dữ liệu mà còn là bảo đảm về tính đúng đắn: nếu ứng dụng có thể ghi vào chỉ mục hoặc dữ liệu khai phá, thì tính toàn vẹn giữa chỉ mục vector và cây tri thức sẽ không còn được bảo đảm.

Dòng \texttt{extra\_\allowbreak{}hosts} khai báo \texttt{host.\allowbreak{}docker.\allowbreak{}internal} trỏ tới cổng của máy chủ. Đây là mấu chốt cho cấu hình mặc định \texttt{QDRANT\_\allowbreak{}URL=\allowbreak{}http:\allowbreak{}/\allowbreak{}/\allowbreak{}host.\allowbreak{}docker.\allowbreak{}internal:\allowbreak{}6333}, cho phép container ứng dụng đọc chỉ mục vector của pipeline đang chạy trên máy chủ mà không cần khởi động thêm một container Qdrant.


\section{Quản lý cấu hình và biến môi trường}
\label{sec:quan-ly-cau-hinh}

\subsection{Kiểm tra cấu hình tại thời điểm khởi động}

Một trong những quyết định kỹ thuật đáng chú ý nhất của backend là việc kiểm tra toàn bộ biến môi trường bằng lược đồ Zod (\texttt{src/\allowbreak{}config/\allowbreak{}env.\allowbreak{}schema.\allowbreak{}ts}) ngay trước khi khởi tạo ứng dụng, thay vì để các giá trị sai lọt xuống tầng nghiệp vụ.

Đoạn khởi động trong \texttt{src/\allowbreak{}main.\allowbreak{}ts} thực hiện kiểm tra này và xử lý lỗi theo cách có chủ ý:

\begin{minted}[fontsize=\footnotesize,breaklines,breakanywhere]{typescript}
try {
  parseEnv(process.env);
} catch (e) {
  if (e instanceof EnvValidationError) {
    console.error(e.message);
    process.exit(1);
  }
  throw e;
}
\end{minted}

Hàm \texttt{parseEnv} dùng \texttt{z.\allowbreak{}object(\{.\allowbreak{}.\allowbreak{}.\allowbreak{}\}).\allowbreak{}safeParse()}, rồi ánh xạ mọi vấn đề thành thông điệp dạng \texttt{TÊN\_\allowbreak{}BIẾN:\allowbreak{} lý do}. Lớp \texttt{EnvValidationError} ghép các thông điệp thành một danh sách, nên khi có nhiều biến sai, toàn bộ danh sách được in ra cùng lúc thay vì phải sửa lại từng lỗi một. Điểm quan trọng về mặt vận hành: tiến trình dừng ngay với mã thoát 1 trước khi kết nối MongoDB, trước khi nạp mô hình và trước khi lắng nghe cổng. Một cấu hình sai bị phát hiện ở thời điểm khởi động rẻ hơn nhiều lần so với việc phát hiện giữa chừng lúc hệ thống đã nhận yêu cầu thật.

\subsection{Biến bắt buộc, biến có giá trị mặc định và ràng buộc độ dài khóa}

Lược đồ phân biệt ba nhóm biến:

\begin{itemize}[label=\(\bullet\)]
    \item \textbf{Bắt buộc} (không có giá trị mặc định): \texttt{MONGO\_\allowbreak{}URI}, \texttt{QDRANT\_\allowbreak{}URL}, \texttt{MODEL\_\allowbreak{}PATH}, \texttt{WORLD\_\allowbreak{}DATA\_\allowbreak{}DIR}, \texttt{MINING\_\allowbreak{}DIR}. Thiếu bất kỳ biến nào cũng khiến tiến trình dừng.
    \item \textbf{Có giá trị mặc định hợp lý}: \texttt{PORT} (3000), \texttt{NODE\_\allowbreak{}ENV} (\texttt{development}), \texttt{QDRANT\_\allowbreak{}COLLECTION} (\texttt{scientific\_\allowbreak{}papers\_\allowbreak{}gold}), \texttt{EMBEDDING\_\allowbreak{}MODEL} (\texttt{nomic-\allowbreak{}ai/\allowbreak{}nomic-\allowbreak{}embed-\allowbreak{}text-\allowbreak{}v1.\allowbreak{}5}), \texttt{LLM\_\allowbreak{}GPU} (\texttt{auto}), \texttt{LLM\_\allowbreak{}CONCURRENCY} (1), \texttt{LLM\_\allowbreak{}MAX\_\allowbreak{}QUEUE} (8), \texttt{LLM\_\allowbreak{}TIMEOUT\_\allowbreak{}MS} (60000), \texttt{CORS\_\allowbreak{}ORIGINS} (\texttt{http:\allowbreak{}/\allowbreak{}/\allowbreak{}localhost:\allowbreak{}5174}).
    \item \textbf{Độ dài tối thiểu của khóa ký}: \texttt{JWT\_\allowbreak{}ACCESS\_\allowbreak{}SECRET} và \texttt{JWT\_\allowbreak{}REFRESH\_\allowbreak{}SECRET} bắt buộc dài ít nhất 32 ký tự, được kiểm tra bằng một hàm trợ giúp \texttt{secret(name)} dùng chung.
\end{itemize}

Việc dùng \texttt{z.\allowbreak{}coerce.\allowbreak{}number().\allowbreak{}int().\allowbreak{}min(min).\allowbreak{}default(def)} cho các biến số cho phép giá trị trong tệp môi trường được ghi dưới dạng chuỗi mà vẫn được kiểm tra đúng kiểu, đồng thời chặn giá trị âm bằng ràng buộc \texttt{min}.

\subsection{Kiểm tra cấu hình riêng cho production}

Lược đồ sử dụng \texttt{superRefine} để áp dụng thêm ba quy tắc chỉ khi \texttt{NODE\_\allowbreak{}ENV =\allowbreak{}=\allowbreak{}=\allowbreak{} 'production'}:

\begin{enumerate}
    \item Cả hai khóa ký đều không được chứa placeholder. Danh sách \texttt{PLACEHOLDER\_\allowbreak{}MARKERS =\allowbreak{} ['change-\allowbreak{}me', 'your-\allowbreak{}secret', 'replace-\allowbreak{}me']} được kiểm tra với phép so khớp không phân biệt hoa thường trên cả hai khóa. Cơ chế này nhắm vào tệp \texttt{.\allowbreak{}env.\allowbreak{}example} đi kèm trong kho mã nguồn \chapterSevenDash{} nếu ai đó sao chép nguyên tệp mẫu sang môi trường thật mà quên thay khóa, ứng dụng sẽ từ chối khởi động.
    \item Hai khóa phải khác nhau. Dùng cùng một giá trị cho access token và refresh token sẽ làm mất sự phân tách giữa hai loại token trong cơ chế token rotation.
    \item \texttt{MONGO\_\allowbreak{}URI} phải chứa thông tin đăng nhập. Quy tắc này áp dụng khi URI bắt đầu bằng \texttt{mongodb:\allowbreak{}/\allowbreak{}/\allowbreak{}} và không có ký tự \texttt{@\allowbreak{}}.
\end{enumerate}

Các quy tắc trên chỉ có hiệu lực ở production. Điều này cần thiết vì ở môi trường phát triển, MongoDB cục bộ không có cơ chế xác thực, và việc bắt buộc có thông tin đăng nhập sẽ làm hỏng quy trình làm việc hằng ngày.

\subsection{Biến môi trường của frontend}

Ở frontend, Vite chỉ thay thế các biến có tiền tố \texttt{VITE\_\allowbreak{}} ngay khi build frontend. Cấu hình chuẩn được đặt trong tệp \texttt{.\allowbreak{}env.\allowbreak{}example}:

\begingroup\small
\begin{longtable}{@{}>{\raggedright\arraybackslash}p{0.31\dimexpr\linewidth-4\tabcolsep\relax}>{\raggedright\arraybackslash}p{0.32\dimexpr\linewidth-4\tabcolsep\relax}>{\raggedright\arraybackslash}p{0.37\dimexpr\linewidth-4\tabcolsep\relax}@{}}
\toprule
\textbf{Biến} & \textbf{Giá trị mặc định} & \textbf{Ý nghĩa} \\
\midrule
\endhead
\texttt{VITE\_\allowbreak{}API\_\allowbreak{}BASE} & \texttt{/\allowbreak{}api/\allowbreak{}v1} & Tiền tố đường dẫn API \\
\texttt{VITE\_\allowbreak{}BACKEND\_\allowbreak{}URL} & \texttt{http:\allowbreak{}/\allowbreak{}/\allowbreak{}localhost:\allowbreak{}3000} & Đích mà proxy của Vite chuyển tiếp trong chế độ phát triển \\
\texttt{VITE\_\allowbreak{}USE\_\allowbreak{}MOCK} & \texttt{true} & \texttt{true} chạy bằng dữ liệu giả lập, \texttt{false} gọi backend thật \\
\texttt{VITE\_\allowbreak{}PIPELINE\_\allowbreak{}BASE} & \texttt{/\allowbreak{}pipeline} & Tiền tố đường dẫn tới pipeline \\
\texttt{VITE\_\allowbreak{}PIPELINE\_\allowbreak{}URL} & \texttt{http:\allowbreak{}/\allowbreak{}/\allowbreak{}localhost:\allowbreak{}8000} & Đích pipeline \\
\texttt{VITE\_\allowbreak{}MOCK\_\allowbreak{}WORLD} & \texttt{true} & \texttt{true} dựng thế giới 3D giả lập, \texttt{false} tải từ pipeline \\
\bottomrule
\end{longtable}
\endgroup

Hai cờ cấu hình \texttt{VITE\_\allowbreak{}USE\_\allowbreak{}MOCK} và \texttt{VITE\_\allowbreak{}MOCK\_\allowbreak{}WORLD} có mặc định là \texttt{true} và được tách bạch theo phạm vi chức năng: một cờ cho luồng hỏi đáp và tra cứu bài báo, cờ còn lại cho thế giới 3D và luồng thu thập dữ liệu. Việc tách bạch này phản ánh thực tế là pipeline chưa phục vụ các tuyến \texttt{/\allowbreak{}world/\allowbreak{}*}, trong khi backend đã phục vụ đầy đủ. Trong \texttt{package.\allowbreak{}json} có hai lệnh tương ứng: \texttt{npm run dev} chạy toàn bộ bằng dữ liệu giả lập và không cần backend, còn \texttt{npm run dev:\allowbreak{}live} đặt \texttt{VITE\_\allowbreak{}USE\_\allowbreak{}MOCK=\allowbreak{}false} để gọi backend thật.

Một hệ quả quan trọng của cơ chế thay thế biến của Vite: giá trị được ghi vào JavaScript bundle khi build. Khi đổi một trong hai cờ cấu hình trong tệp \texttt{.\allowbreak{}env}, phải build lại frontend, nếu không trình duyệt vẫn dùng giá trị cũ. Với bản triển khai dạng container image, điều này có nghĩa là phải build lại image chứ không chỉ khởi động lại container.

\subsection{Phạm vi sử dụng secret ở backend và frontend}

Khóa ký JWT, khóa truy cập Qdrant và thông tin đăng nhập MongoDB chỉ thuộc về phía backend. Giao diện web được phục vụ dưới dạng tệp tĩnh và không chứa secret; mọi yêu cầu có trạng thái đều đi qua API với header \texttt{Authorization:\allowbreak{} Bearer}. Điều này đồng nghĩa với việc tệp \texttt{.\allowbreak{}env} của frontend không được chứa secret, và các biến \texttt{VITE\_\allowbreak{}*} \chapterSevenDash{} vốn được nhúng công khai vào gói JavaScript \chapterSevenDash{} tuyệt đối không được dùng để lưu thông tin xác thực.


\section{Triển khai Backend, Frontend và Database}
\label{sec:trien-khai-be-fe-db}

\subsection{Bước 1 \texorpdfstring{{\normalfont\fontfamily{lmr}\selectfont\textemdash}}{—} Chuẩn bị chỉ mục truy hồi}

Theo mặc định, ứng dụng trong container truy cập Qdrant của pipeline qua \texttt{host.\allowbreak{}docker.\allowbreak{}internal:\allowbreak{}6333} và đọc collection \texttt{scientific\_\allowbreak{}papers\_\allowbreak{}gold}. Trước khi thử tìm kiếm, cần xác nhận chỉ mục đã được nạp và chứa đúng số chiều vector.

Cơ chế kiểm tra số chiều được cài đặt trong \texttt{VectorStoreService.\allowbreak{}onModuleInit()}: khi khởi động, service đọc cấu hình collection và so sánh kích thước vector với hằng số \texttt{EMBEDDING\_\allowbreak{}DIM} (768 chiều). Nếu lệch, service ghi cảnh báo \texttt{QDRANT\_\allowbreak{}DIMENSION\_\allowbreak{}MISMATCH}. Cùng lúc đó, service kiểm tra sự tồn tại của trường \texttt{twig\_\allowbreak{}id} trong lược đồ payload; nếu thiếu, nó ghi cảnh báo \texttt{QDRANT\_\allowbreak{}NO\_\allowbreak{}TWIG\_\allowbreak{}ID} và nội dung cảnh báo cho biết chức năng so sánh sẽ không được giới hạn theo nhánh.

Hai cảnh báo này là ví dụ điển hình của triết lý "lỗi cục bộ, không làm sập hệ thống" đã được Chương 3 xác lập. Sai cấu hình chỉ mục không làm ứng dụng dừng, nhưng luôn để lại dấu vết trong nhật ký để người vận hành phát hiện.

Khi chạy độc lập không có pipeline, bản hồ Compose \texttt{standalone} tạo một Qdrant rỗng. Trường hợp này cần đặt \texttt{RAGTS\_\allowbreak{}QDRANT\_\allowbreak{}URL=\allowbreak{}http:\allowbreak{}/\allowbreak{}/\allowbreak{}qdrant:\allowbreak{}6333}, nạp chỉ mục bằng \texttt{scripts/\allowbreak{}seed-\allowbreak{}qdrant.\allowbreak{}ts} và chọn đúng collection trước khi thử tìm kiếm. Tệp này có một biện pháp an toàn đáng chú ý: nó từ chối mọi collection không có hậu tố \texttt{\_\allowbreak{}test}, nhằm ngăn việc vô tình ghi đè chỉ mục thật.

\subsection{Bước 2 \texorpdfstring{{\normalfont\fontfamily{lmr}\selectfont\textemdash}}{—} Khởi động MongoDB và API}

Sau khi chuẩn bị mô hình và cấp các biến môi trường, chạy lệnh sau từ thư mục \texttt{backend\_\allowbreak{}rag}:

\begin{minted}[fontsize=\footnotesize,breaklines,breakanywhere]{bash}
docker compose up -d --build
\end{minted}

Trước khi chạy, cần cấp tối thiểu hai khóa ký cùng khóa truy cập Qdrant:

\begin{minted}[fontsize=\footnotesize,breaklines,breakanywhere]{bash}
export JWT_ACCESS_SECRET=...
export JWT_REFRESH_SECRET=...
export QDRANT_API_KEY=...
\end{minted}

Tệp Compose sử dụng cú pháp \texttt{\$\{JWT\_\allowbreak{}ACCESS\_\allowbreak{}SECRET:\allowbreak{}?set JWT\_\allowbreak{}ACCESS\_\allowbreak{}SECRET (>=\allowbreak{} 32 chars)\}}. Toán tử \texttt{:\allowbreak{}?} khiến Compose dừng lại với thông báo lỗi nếu biến đó không được đặt hoặc rỗng. Đây là cơ chế kiểm tra cấu hình tại tầng Compose, bổ sung cho lớp kiểm tra Zod ở tầng ứng dụng. Hai lớp bảo vệ có phạm vi khác nhau: Compose ngăn khởi động container với cấu hình thiếu, còn Zod ngăn ứng dụng chạy được với cấu hình sai nghĩa.

Về mặt mạng, cả MongoDB (cổng 27018) và API (cổng 3000) mặc định chỉ mở trên địa chỉ loopback của máy chủ. Điều này đúng với mô hình triển khai đề xuất, trong đó Nginx là thành phần duy nhất tiếp nhận kết nối từ bên ngoài.

\subsection{Bước 3 \texorpdfstring{{\normalfont\fontfamily{lmr}\selectfont\textemdash}}{—} Đưa giao diện vào hoạt động}

Build Docker image từ thư mục \texttt{frontend\_\allowbreak{}rag}, cung cấp cấu hình Nginx và nối container frontend vào cùng mạng với dịch vụ \texttt{app}. Chuỗi kiểm tra sau khi triển khai nên diễn ra theo thứ tự từ lớp trong ra ngoài:

\begin{enumerate}
    \item Kiểm tra liveness của container: \texttt{GET /\allowbreak{}api/\allowbreak{}v1/\allowbreak{}health/\allowbreak{}live} trả HTTP 200.
    \item Kiểm tra readiness của backend: \texttt{GET /\allowbreak{}api/\allowbreak{}v1/\allowbreak{}health/\allowbreak{}ready}.
    \item Kiểm tra qua địa chỉ công khai: mở một đường dẫn trực tiếp của giao diện, xác nhận trang dựng đầy đủ.
    \item Kiểm tra chức năng nghiệp vụ: gửi một câu hỏi có bài báo đã biết trong chỉ mục, quan sát trình duyệt nhận dần câu trả lời và mở được đúng bài báo từ nguồn trích dẫn.
\end{enumerate}

Bước cuối cùng là bước có giá trị chứng nhận cao nhất, vì nó kiểm tra đồng thời toàn bộ quy trình xử lý yêu cầu: nhúng câu hỏi, định tuyến trên cây tri thức, truy vấn Qdrant, xếp hạng lại, sinh văn bản và truyền theo luồng SSE. Ở bản production, bước này phải thực hiện qua địa chỉ HTTPS chứ không phải qua địa chỉ cục bộ, để đồng thời xác nhận chứng thư, chuyển hướng và luồng SSE đi qua proxy.

Vì Nginx trong dự án được vận hành bằng hai tệp Compose riêng ở thư mục gốc kho, trình tự triển khai đầy đủ cho bản production là:

\begin{minted}[fontsize=\footnotesize,breaklines,breakanywhere]{bash}
# 1. Đặt chứng thư và khóa riêng (thư mục này đã nằm trong .gitignore)
#    ssl/cert.pem và ssl/key.pem

# 2. Dựng bản phân phối giao diện — Nginx ánh xạ thư mục này vào /usr/share/nginx/html
cd frontend_rag && npm run build && cd ..

# 3. Khởi động toàn bộ stack
docker compose -f docker-compose.prod.yaml up -d --build
\end{minted}

Chi tiết về hai tệp Compose và cách ánh xạ cấu hình Nginx được trình bày ở Mục 7.6.2.

\subsection{Các thành phần, cổng kết nối và cách kiểm tra}

\begingroup\small
\begin{longtable}{@{}>{\raggedright\arraybackslash}p{0.22\dimexpr\linewidth-6\tabcolsep\relax}>{\raggedright\arraybackslash}p{0.18\dimexpr\linewidth-6\tabcolsep\relax}>{\raggedright\arraybackslash}p{0.31\dimexpr\linewidth-6\tabcolsep\relax}>{\raggedright\arraybackslash}p{0.29\dimexpr\linewidth-6\tabcolsep\relax}@{}}
\toprule
\textbf{Thành phần} & \textbf{Cổng (máy chủ)} & \textbf{Nguồn cấu hình} & \textbf{Kiểm tra} \\
\midrule
\endhead
MongoDB 7 (replica set \texttt{rs0}) & 27018 & \texttt{backend\_\allowbreak{}rag/\allowbreak{}docker-\allowbreak{}compose.\allowbreak{}yml} & \texttt{rs.\allowbreak{}status().\allowbreak{}ok} \\
API NestJS & 3000 & \texttt{backend\_\allowbreak{}rag/\allowbreak{}docker-\allowbreak{}compose.\allowbreak{}yml} & \texttt{/\allowbreak{}api/\allowbreak{}v1/\allowbreak{}health/\allowbreak{}live} \\
Nginx (phát triển) & 80 & \texttt{docker-\allowbreak{}compose.\allowbreak{}yaml} + \texttt{nginx.\allowbreak{}dev.\allowbreak{}conf} & Tải trang chủ \\
Nginx (production) & 80 \ensuremath{\rightarrow} 301, 443 & \texttt{docker-\allowbreak{}compose.\allowbreak{}prod.\allowbreak{}yaml} + \texttt{nginx.\allowbreak{}prod.\allowbreak{}conf} & Kiểm tra chuỗi tin cậy \\
Giao diện tĩnh & \chapterSevenDash{} & \texttt{frontend\_\allowbreak{}rag/\allowbreak{}dist} \ensuremath{\rightarrow} \texttt{/\allowbreak{}usr/\allowbreak{}share/\allowbreak{}nginx/\allowbreak{}html} & Tải trang chủ \\
Qdrant (chỉ đọc) & 6333 (ngoài container) & \texttt{RAGTS\_\allowbreak{}QDRANT\_\allowbreak{}URL} & \texttt{getCollection} \\
Mô hình ngôn ngữ & \chapterSevenDash{} & \texttt{/\allowbreak{}models} (chỉ đọc) & Sự kiện \texttt{LLM\_\allowbreak{}READY} \\
Mô hình ngôn ngữ & \chapterSevenDash{} & \texttt{/\allowbreak{}models} (chỉ đọc) & Sự kiện \texttt{LLM\_\allowbreak{}READY} \\
\bottomrule
\end{longtable}
\endgroup


\section{Triển khai Local AI Runtime và các model cần thiết}
\label{sec:trien-khai-local-ai}

\subsection{Hai mô hình đảm nhận hai bước khác nhau}

Cấu hình suy luận cục bộ của RAG\_TS dựa trên hai mô hình độc lập, mỗi mô hình phục vụ một bước khác nhau trong đường xử lý.

Transformers.js nạp \textbf{Nomic Embed Text v1.5} để ánh xạ câu hỏi vào cùng không gian 768 chiều với các đoạn bài báo trong Qdrant. Mô hình embedding này được chỉ định qua \texttt{EMBEDDING\_\allowbreak{}MODEL}, mặc định là \texttt{nomic-\allowbreak{}ai/\allowbreak{}nomic-\allowbreak{}embed-\allowbreak{}text-\allowbreak{}v1.\allowbreak{}5}, và có thể trỏ tới một mã trên Hugging Face hoặc một thư mục cục bộ.

Sau bước truy hồi, \texttt{node-\allowbreak{}llama-\allowbreak{}cpp} chạy \textbf{Qwen2.5-7B-Instruct}, bản lượng tử hóa Q4\_K\_M. Mô hình được trỏ tới qua \texttt{MODEL\_\allowbreak{}PATH}, mặc định trong Compose là \texttt{/\allowbreak{}models/\allowbreak{}\$\{RAGTS\_\allowbreak{}MODEL\_\allowbreak{}FILE:\allowbreak{}-\allowbreak{}qwen2.\allowbreak{}5-\allowbreak{}7b-\allowbreak{}instruct-\allowbreak{}q4\_\allowbreak{}k\_\allowbreak{}m.\allowbreak{}gguf\}}.

Cảnh báo vận hành quan trọng nhất của phần này: \textbf{thay mô hình embedding đòi hỏi phải lập chỉ mục lại bằng chính mô hình mới}. Nguyên nhân nằm ở việc hai mô hình khác nhau tạo ra hai hệ toạ độ vector khác nhau, kể cả khi số chiều trùng khớp. Khi đó, vector câu hỏi sẽ được so sánh với các vector đoạn văn bản được tạo bởi một mô hình hoàn toàn khác, và kết quả truy hồi sẽ là ngẫu nhiên theo đúng nghĩa đen. Đây là một trong những lỗi cấu hình đắt giá nhất mà dấu hiệu bên ngoài rất dễ gây hiểu nhầm: hệ thống vẫn trả lời, vẫn trích dẫn, nhưng chất lượng truy hồi suy giảm âm thầm. Cơ chế kiểm tra \texttt{QDRANT\_\allowbreak{}DIMENSION\_\allowbreak{}MISMATCH} chỉ bắt được trường hợp sai \emph{số chiều}, không bắt được trường hợp sai \emph{mô hình cùng số chiều}.

\subsection{Tải trọng số}

Trên máy chuẩn bị mô hình, có thể chạy lệnh sau từ \texttt{backend\_\allowbreak{}rag} sau khi cài các phụ thuộc phát triển:

\begin{minted}[fontsize=\footnotesize,breaklines,breakanywhere]{bash}
npm run download:model
\end{minted}

Tệp \texttt{scripts/\allowbreak{}download-\allowbreak{}model.\allowbreak{}ts} tải từ địa chỉ \texttt{https:\allowbreak{}/\allowbreak{}/\allowbreak{}huggingface.\allowbreak{}co/\allowbreak{}bartowski/\allowbreak{}Qwen2.\allowbreak{}5-\allowbreak{}7B-\allowbreak{}Instruct-\allowbreak{}GGUF/\allowbreak{}resolve/\allowbreak{}main/\allowbreak{}Qwen2.\allowbreak{}5-\allowbreak{}7B-\allowbreak{}Instruct-\allowbreak{}Q4\_\allowbreak{}K\_\allowbreak{}M.\allowbreak{}gguf} với dung lượng khoảng 4,68 GB. Script có hai chi tiết vận hành đáng lưu ý:

\begin{itemize}[label=\(\bullet\)]
    \item Thứ tự ưu tiên đường dẫn đích là \texttt{MODEL\_\allowbreak{}PATH}, rồi \texttt{LLM\_\allowbreak{}MODEL\_\allowbreak{}PATH}, rồi mặc định \texttt{.\allowbreak{}.\allowbreak{}/\allowbreak{}model/\allowbreak{}qwen2.\allowbreak{}5-\allowbreak{}7b-\allowbreak{}instruct-\allowbreak{}q4\_\allowbreak{}k\_\allowbreak{}m.\allowbreak{}gguf}. Nhờ vậy, cùng một lệnh tải có thể dùng được cho cả chạy ngoài Docker lẫn chạy trong container.
    \item Script kiểm tra kích thước tệp đã tồn tại: nếu lớn hơn 4 GB thì coi như tải xong và bỏ qua, còn nếu nhỏ hơn thì coi như bản tải dở dang và tải lại từ đầu. Ngưỡng 4 GB được chọn dựa trên việc dung lượng thực tế của bản Q4\_K\_M là 4,68 GB, nên một tệp dưới ngưỡng này chắc chắn là bản không hoàn chỉnh. Script cũng in tiến độ theo từng mốc 50 MB.
\end{itemize}

Mô hình embedding cũng cần truy cập nguồn tải ở lần nạp đầu tiên, trừ khi đã có sẵn mô hình cục bộ trong bộ nhớ đệm \texttt{hf-\allowbreak{}cache}.

\subsection{Nạp mô hình ở chế độ nền và xử lý yêu cầu trong thời gian khởi động}

Mô hình 7B mất một khoảng thời gian đáng kể để nạp, nên \texttt{LlmService} nạp ở chế độ nền thay vì chặn quá trình khởi động:

\begin{minted}[fontsize=\footnotesize,breaklines,breakanywhere]{typescript}
onModuleInit(): void {
  // Loading a 7B model takes a while: do it in the background, health reports `ready`.
  void this.load();
}
\end{minted}

Trong \texttt{onModuleInit}, \texttt{void this.load()} khởi chạy tác vụ nạp mô hình mà không chặn quá trình khởi động của API. Trong lúc mô hình đang nạp, cờ \texttt{ready} vẫn là \texttt{false}. Endpoint hỏi đáp trả HTTP~503 kèm mã \texttt{llm unavailable}, giúp người vận hành phân biệt lỗi mô hình chưa sẵn sàng với lỗi truy cập chỉ mục hoặc cơ sở dữ liệu.

Tài liệu vận hành của dự án ghi nhận thời gian nạp khoảng 5 giây trên nền Metal. Cần lưu ý rằng Docker trên macOS không có Metal, nên LLM sẽ chạy trên CPU.

\subsection{Cơ chế hàng đợi và giới hạn tải}

\texttt{LlmService} bọc engine sinh văn bản trong một hàng đợi giới hạn (\texttt{LlmQueue}) với hai tham số: số lượt sinh song song (\texttt{LLM\_\allowbreak{}CONCURRENCY}, mặc định 1) và độ dài hàng chờ tối đa (\texttt{LLM\_\allowbreak{}MAX\_\allowbreak{}QUEUE}, mặc định 8).

Cơ chế hàng đợi có hai đặc tính thiết kế đáng chú ý:

Thứ nhất, khi hàng đợi đầy, hệ thống từ chối có kiểm soát bằng mã \texttt{429 Too Many Requests} thay vì để yêu cầu chờ vô hạn. Đây là cơ chế giảm tải khi bão hòa: từ chối sớm bảo toàn tài nguyên cho những người dùng đang được phục vụ, thay vì để toàn bộ người dùng cùng chờ.

Thứ hai, khi một người dùng ngắt kết nối giữa chừng, lượt chờ của họ rời khỏi hàng đợi ngay lập tức. Cơ chế này dựa trên \texttt{AbortSignal}: mỗi lượt chờ đăng ký một trình xử lý huỷ, khi tín hiệu huỷ phát ra, lượt chờ tương ứng được rút khỏi mảng và bị từ chối. Điều này tránh tình huống một người dùng đã đóng trình duyệt vẫn chiếm chỗ trong hàng đợi.

Ngoài ra, \texttt{LLM\_\allowbreak{}TIMEOUT\_\allowbreak{}MS} đặt timeout cho một lượt sinh. Khi hết thời gian, yêu cầu bị hủy và trả về lỗi \texttt{504} với mã \texttt{llm timeout}. Lỗi \texttt{503} với mã \texttt{llm unavailable} được dùng cho trường hợp mô hình chưa nạp xong hoặc gặp lỗi sinh.

Sự phân biệt này quan trọng cho vận hành: \texttt{503} cho biết hệ thống chưa sẵn sàng (có thể tự phục hồi), còn \texttt{429} cho biết hệ thống đã sẵn sàng nhưng đang quá tải (cần giảm tải hoặc tăng tài nguyên), còn \texttt{504} cho biết một lượt cụ thể vượt thời gian cho phép.

\subsection{Cấu hình hàng đợi trong Compose}

Compose đặt mặc định \texttt{LLM\_\allowbreak{}GPU=\allowbreak{}auto}, \texttt{LLM\_\allowbreak{}MAX\_\allowbreak{}QUEUE=\allowbreak{}8} và \texttt{LLM\_\allowbreak{}TIMEOUT\_\allowbreak{}MS=\allowbreak{}120000}. Tệp môi trường mẫu \texttt{.\allowbreak{}env.\allowbreak{}example} đặt timeout là 60 giây. Sự khác biệt này là có chủ ý: giá trị trong Compose là giá trị cho môi trường thật, nơi phần cứng mạnh hơn và thời gian sinh có thể dài hơn; còn giá trị trong tệp mẫu là giá trị an toàn cho máy phát triển.

Khi đánh giá độ trễ, cần ghi lại đồng thời cấu hình thực thi, thiết bị mà runtime chọn, và thời gian chờ trong hàng đợi. Ba thông tin này phân biệt được hai nguyên nhân hoàn toàn khác nhau: nghẽn do nhiều người dùng đồng thời, và nghẽn do bản thân tốc độ sinh token của mô hình.

\subsection{Ghi chú về gói Laya và xếp hạng lại}

Ngoài nhúng và sinh văn bản, backend hỗ trợ gói ONNX Laya dùng để xếp hạng lại các đoạn truy hồi, đánh giá bằng chứng và kiểm tra mức độ câu trả lời bám sát nguồn. Khi sử dụng, người vận hành gắn thư mục mô hình và truyền đường dẫn trong container qua biến \texttt{LAYA\_\allowbreak{}MODEL\_\allowbreak{}DIR}. Khi để trống đường dẫn, backend bỏ qua các bước dùng Laya \chapterSevenDash{} cơ chế tắt tính năng theo kiểu "không cấu hình thì không chạy", tương tự cách \texttt{WORLD\_\allowbreak{}DATA\_\allowbreak{}DIR} và \texttt{MINING\_\allowbreak{}DIR} xử lý.


\section{Cấu hình Nginx và Reverse Proxy}
\label{sec:cau-hinh-nginx}

\subsection{Hai bộ cấu hình Nginx trong dự án}

Hệ thống có hai bộ cấu hình Nginx riêng biệt, cùng nằm ở thư mục gốc kho mã nguồn và được thêm vào trong commit \texttt{3580935} với thông điệp \emph{"feat: setup NginX for develop and production environment with Docker"}. Việc tách hai bộ cấu hình nhằm đáp ứng hai yêu cầu vận hành vốn không thể dung hòa trong một tệp duy nhất: bản phát triển phải chạy được ngay với HTTP thuần trên \texttt{localhost}, còn bản production phải thực hiện TLS termination tại Nginx và chuyển hướng toàn bộ lưu lượng HTTP sang HTTPS.

\begingroup\small
\begin{longtable}{@{}>{\raggedright\arraybackslash}p{0.22\dimexpr\linewidth-6\tabcolsep\relax}>{\raggedright\arraybackslash}p{0.18\dimexpr\linewidth-6\tabcolsep\relax}>{\raggedright\arraybackslash}p{0.31\dimexpr\linewidth-6\tabcolsep\relax}>{\raggedright\arraybackslash}p{0.29\dimexpr\linewidth-6\tabcolsep\relax}@{}}
\toprule
\textbf{Tệp} & \textbf{Mục đích} & \textbf{Cổng} & \textbf{Điểm khác biệt chính} \\
\midrule
\endhead
\texttt{nginx.\allowbreak{}dev.\allowbreak{}conf} & Môi trường phát triển & 80 & Không TLS, không chuyển hướng, không log riêng \\
\texttt{nginx.\allowbreak{}prod.\allowbreak{}conf} & Môi trường production & 80 \ensuremath{\rightarrow} 301, 443 & TLS termination, chuyển hướng HTTPS, log riêng, hỗ trợ WebSocket \\
\bottomrule
\end{longtable}
\endgroup

Cấu hình production, lấy từ \texttt{nginx.\allowbreak{}prod.\allowbreak{}conf}:

\begin{minted}[fontsize=\footnotesize,breaklines,breakanywhere]{nginx}
events {
    worker_connections 1024;
}

http {
    include mime.types;

    upstream backend_rag_cluster {
        server backend-rag:8000 max_fails=3 fail_timeout=30s; # 'backend-rag' là tên service trong Docker
    }

    server {
        listen 80;
        server_name tenmien.com www.tenmien.com;
        return 301 https://$host$request_uri;
    }

    server {
        listen 443 ssl;
        server_name tenmien.com www.tenmien.com;

        ssl_certificate "/etc/nginx/ssl/cert.pem";
        ssl_certificate_key "/etc/nginx/ssl/key.pem";

        access_log /var/log/nginx/rag_access.log;
        error_log /var/log/nginx/rag_error.log warn;

        proxy_read_timeout 300;
        proxy_connect_timeout 300;
        proxy_send_timeout 300;

        location / {
            root "/usr/share/nginx/html";
            try_files $uri /index.html;
        }

        location /api/ {
            proxy_pass http://backend_rag_cluster;

            proxy_redirect off;
            proxy_buffering off;

            proxy_next_upstream error timeout invalid_header http_500 http_502 http_503 http_504;

            proxy_set_header Host $host;
            proxy_set_header X-Real-IP $remote_addr;
            proxy_set_header X-Forwarded-For $proxy_add_x_forwarded_for;
            proxy_set_header X-Forwarded-Proto $scheme;

            # Hỗ trợ WebSockets
            proxy_set_header Upgrade $http_upgrade;
            proxy_set_header Connection "upgrade";
        }
    }
}
\end{minted}

Cấu hình phát triển, lấy từ \texttt{nginx.\allowbreak{}dev.\allowbreak{}conf}:

\begin{minted}[fontsize=\footnotesize,breaklines,breakanywhere]{nginx}
events {}

http {
    include mime.types;

    upstream backend_rag_cluster {
        server backend-rag:8000;
    }

    server {
        listen 80;
        server_name localhost;

        location / {
            root /usr/share/nginx/html;
            try_files $uri /index.html;
        }

        location /api/ {
            proxy_pass http://backend_rag_cluster;
            proxy_set_header Host $host;
            proxy_set_header X-Real-IP $remote_addr;
        }
    }
}
\end{minted}

Bảy quyết định thiết kế đáng chú ý trong hai bộ cấu hình này:

\textbf{Thứ nhất, sử dụng khối \texttt{upstream} thay vì trỏ thẳng vào địa chỉ container.} Khối \texttt{upstream backend\_\allowbreak{}rag\_\allowbreak{}cluster} định nghĩa một tên logic cho backend, tách cấu hình định tuyến khỏi tên service thực tế trong Docker. Điều này cho phép chuyển backend sang nhiều bản sao chỉ bằng cách sửa một khối, mà không phải sửa mọi khối \texttt{location}. Ở bản production, khối này còn gắn hai tham số \texttt{max\_\allowbreak{}fails=\allowbreak{}3} và \texttt{fail\_\allowbreak{}timeout=\allowbreak{}30s}, quy định rằng Nginx sẽ đánh dấu một bản sao là lỗi sau 3 lần thất bại và chỉ thử lại sau 30 giây.

\textbf{Thứ hai, \texttt{proxy\_\allowbreak{}buffering off} xuất hiện ở bản production nhưng không có ở bản phát triển.} Đây là điểm cốt lõi cho luồng SSE. Khi Nginx bật response buffering theo mặc định, các phần nhỏ sẽ bị gom lại trước khi chuyển tiếp, khiến giao diện hiển thị câu trả lời theo từng đoạn lớn thay vì theo dòng, và làm chậm thời điểm người dùng nhìn thấy token đầu tiên. Bật \texttt{proxy\_\allowbreak{}buffering off} ở bản production nhưng bỏ ở bản phát triển phản ánh đúng đặc thù: chỉ bản triển khai thật mới cần phục vụ luồng SSE, còn bản phát triển chỉ phục vụ giao diện và các lời gọi API không phát trực tiếp.

\textbf{Thứ ba, ba giá trị thời gian chờ được đặt cùng nhau ở mức 300 giây.} \texttt{proxy\_\allowbreak{}read\_\allowbreak{}timeout}, \texttt{proxy\_\allowbreak{}connect\_\allowbreak{}timeout} và \texttt{proxy\_\allowbreak{}send\_\allowbreak{}timeout} đều bằng 300, thay vì mặc định 60 giây của Nginx. Con số này phải lớn hơn timeout suy luận của LLM. Đây là mối liên hệ trực tiếp với cấu hình ở Mục 7.5: nếu \texttt{LLM\_\allowbreak{}TIMEOUT\_\allowbreak{}MS} được đặt là 120000 (120 giây) trong Compose, thì thời gian chờ 300 giây của Nginx tạo biên an toàn hợp lý \chapterSevenDash{} Nginx chờ lâu hơn backend, nên khi backend hết thời gian và trả lỗi \texttt{504}, người dùng nhận được thông báo lỗi có nghĩa từ ứng dụng thay vì một trang lỗi 504 trống từ proxy. Ngược lại, nếu đặt thời gian chờ của Nginx ngắn hơn timeout của LLM, người dùng sẽ thấy lỗi proxy trong khi mô hình vẫn đang sinh câu trả lời, tức là tài nguyên GPU bị lãng phí.

\textbf{Thứ tư, Nginx chuyển tiếp bốn HTTP header tới backend.} Ba header đầu (\texttt{Host}, \texttt{X-\allowbreak{}Real-\allowbreak{}IP}, \texttt{X-\allowbreak{}Forwarded-\allowbreak{}For}) là bắt buộc để backend nhận được địa chỉ client thật, và \texttt{X-\allowbreak{}Forwarded-\allowbreak{}Proto} cho biết giao thức gốc là gì. Điều này gắn trực tiếp với cơ chế giới hạn tần suất được trình bày ở Mục 7.6.3: nếu Nginx không gửi \texttt{X-\allowbreak{}Forwarded-\allowbreak{}For}, Express sẽ coi mọi yêu cầu là đến từ một địa chỉ duy nhất là Nginx, và toàn bộ người dùng sẽ dùng chung một bộ đếm.

\textbf{Thứ năm, hỗ trợ WebSocket được cấu hình sẵn.} Hai dòng \texttt{proxy\_\allowbreak{}set\_\allowbreak{}header Upgrade \$http\_\allowbreak{}upgrade;} và \texttt{proxy\_\allowbreak{}set\_\allowbreak{}header Connection "upgrade";} cho phép các kết nối WebSocket đi xuyên qua proxy. Hiện tại hệ thống chưa dùng WebSocket \chapterSevenDash{} luồng SSE dùng HTTP chunked \chapterSevenDash{} nhưng việc cấu hình sẵn giúp không phải sửa cấu hình reverse proxy khi bổ sung một kênh truyền thời gian thực khác.

\textbf{Thứ sáu, \texttt{proxy\_\allowbreak{}next\_\allowbreak{}upstream} liệt kê đầy đủ các mã lỗi.} Cấu hình khai báo \texttt{error timeout invalid\_\allowbreak{}header http\_\allowbreak{}500 http\_\allowbreak{}502 http\_\allowbreak{}503 http\_\allowbreak{}504}, nghĩa là Nginx sẽ thử bản sao backend tiếp theo khi bản sao hiện tại trả về một trong các lỗi đó. Cơ chế failover chỉ hoạt động khi khối \texttt{upstream} khai báo nhiều hơn một bản sao; với cấu hình hiện tại chỉ có một \texttt{server}, Nginx không có backend thứ hai để thử lại.

\textbf{Thứ bảy, nhật ký được tách thành hai tệp riêng.} \texttt{access\_\allowbreak{}log /\allowbreak{}var/\allowbreak{}log/\allowbreak{}nginx/\allowbreak{}rag\_\allowbreak{}access.\allowbreak{}log} ghi lại các yêu cầu, còn \texttt{error\_\allowbreak{}log /\allowbreak{}var/\allowbreak{}log/\allowbreak{}nginx/\allowbreak{}rag\_\allowbreak{}error.\allowbreak{}log warn} chỉ ghi các sự kiện mức cảnh báo trở lên. Hai tệp này được ánh xạ vào volume \texttt{prod-\allowbreak{}logs} (bản production) và \texttt{dev-\allowbreak{}logs} (bản phát triển), nên nhật ký của Nginx tồn tại lâu hơn vòng đời container và có thể thu thập bằng \texttt{docker logs} hoặc đẩy sang hệ thống tổng hợp.

\subsection{Đưa Nginx vào vận hành: hai tệp Compose}

Nginx không được khởi động trong \texttt{backend\_\allowbreak{}rag/\allowbreak{}docker-\allowbreak{}compose.\allowbreak{}yml} mà bằng hai tệp Compose riêng ở thư mục gốc kho, phân biệt rõ môi trường.

Tệp \texttt{docker-\allowbreak{}compose.\allowbreak{}yaml} cho môi trường phát triển:

\begin{minted}[fontsize=\footnotesize,breaklines,breakanywhere]{yaml}
services:
  backend-rag:
    build: ./backend-rag
    ports:
      - "8000:8000"
    environment:
      - QDRANT_URL=
      - QDRANT_API_KEY=
      - MINIO_ENDPOINT=
      - MINIO_ACCESS_KEY=
      - MINIO_SECRET_KEY=

  # Container cho NginX (Web Server & Reverse Proxy)
  nginx:
    image: nginx:alpine
    ports:
      - "80:80"
    volumes:
      # Ánh xạ file cấu hình dành riêng cho dev
      - ./nginx.dev.conf:/etc/nginx/nginx.conf
      - ./frontend-rag/dist:/usr/share/nginx/html
      - ./dev-logs:/var/log/nginx
    depends_on:
      - backend-rag
\end{minted}

Tệp \texttt{docker-\allowbreak{}compose.\allowbreak{}prod.\allowbreak{}yaml} cho môi trường production:

\begin{minted}[fontsize=\footnotesize,breaklines,breakanywhere]{yaml}
services:
  backend-rag:
    build: ./backend-rag
    ports:
      - "8000:8000"
    environment:
      - QDRANT_URL=
      - QDRANT_API_KEY=
      - MINIO_ENDPOINT=
      - MINIO_ACCESS_KEY=
      - MINIO_SECRET_KEY=
  nginx:
    image: nginx:alpine
    ports:
      - "80:80"
      - "443:443"
    volumes:
      # Ánh xạ file cấu hình production (có chứa tên miền thật)
      - ./nginx.prod.conf:/etc/nginx/nginx.conf
      - ./ssl:/etc/nginx/ssl
      - ./frontend-rag/dist:/usr/share/nginx/html
      - ./prod-logs:/var/log/nginx
    depends_on:
      - backend-rag
\end{minted}

Bốn điểm khác biệt giữa hai bản cần lưu ý khi vận hành:

\begingroup\small
\begin{longtable}{@{}>{\raggedright\arraybackslash}p{0.31\dimexpr\linewidth-4\tabcolsep\relax}>{\raggedright\arraybackslash}p{0.32\dimexpr\linewidth-4\tabcolsep\relax}>{\raggedright\arraybackslash}p{0.37\dimexpr\linewidth-4\tabcolsep\relax}@{}}
\toprule
\textbf{Hạng mục} & \textbf{Phát triển} & \textbf{Production} \\
\midrule
\endhead
Cổng công khai & 80 & 80 và 443 \\
Tệp cấu hình & \texttt{nginx.\allowbreak{}dev.\allowbreak{}conf} & \texttt{nginx.\allowbreak{}prod.\allowbreak{}conf} \\
Gắn thư mục chứng thư & Không & \texttt{.\allowbreak{}/\allowbreak{}ssl:\allowbreak{}/\allowbreak{}etc/\allowbreak{}nginx/\allowbreak{}ssl} \\
Volume nhật ký & \texttt{dev-\allowbreak{}logs} & \texttt{prod-\allowbreak{}logs} \\
\bottomrule
\end{longtable}
\endgroup

Điểm thứ nhất: gắn thư mục \texttt{.\allowbreak{}/\allowbreak{}ssl} vào \texttt{/\allowbreak{}etc/\allowbreak{}nginx/\allowbreak{}ssl} chỉ xuất hiện ở bản production. Thư mục \texttt{ssl/\allowbreak{}} chứa \texttt{cert.\allowbreak{}pem} và \texttt{key.\allowbreak{}pem} được thêm vào \texttt{.\allowbreak{}gitignore} ngay từ commit đầu tiên, nên tệp chứng thư và khóa riêng không bao giờ được commit lên kho mã nguồn \chapterSevenDash{} đây là yêu cầu bắt buộc, không phải lựa chọn. Người vận hành phải tự đặt hai tệp này vào thư mục trước khi khởi động bản production.

Điểm thứ hai: cả hai bản đều ánh xạ thư mục \texttt{frontend-\allowbreak{}rag/\allowbreak{}dist} vào \texttt{/\allowbreak{}usr/\allowbreak{}share/\allowbreak{}nginx/\allowbreak{}html} thay vì sao chép tệp vào ảnh. Cách làm này cho phép triển khai lại giao diện bằng cách build lại frontend rồi tải lại container Nginx, mà không cần build lại image.

Điểm thứ ba: khối \texttt{depends\_\allowbreak{}on:\allowbreak{} -\allowbreak{} backend-\allowbreak{}rag} chỉ bảo đảm thứ tự khởi động chứ không bảo đảm backend đã sẵn sàng nhận yêu cầu. Với Nginx, tác động của điều này ít nghiêm trọng hơn so với backend vì các yêu cầu đến sớm sẽ chỉ nhận lỗi 502 tạm thời từ \texttt{proxy\_\allowbreak{}next\_\allowbreak{}upstream}, sau đó được phục vụ bình thường khi backend sẵn sàng.

Điểm thứ tư: tên service \texttt{backend-\allowbreak{}rag} trong Compose khớp với tên trong khối \texttt{upstream} của cả hai tệp cấu hình Nginx. Đây là điểm liên hệ bắt buộc phải giữ đồng bộ: nếu đổi tên service trong Compose mà không sửa \texttt{nginx.\allowbreak{}dev.\allowbreak{}conf} và \texttt{nginx.\allowbreak{}prod.\allowbreak{}conf}, Nginx sẽ không phân giải được tên và mọi yêu cầu API sẽ trả lỗi 502.

\subsection{Cấu hình \texttt{trust proxy} và rate limiting theo địa chỉ IP}

Khi dùng địa chỉ IP máy khách để giới hạn tần suất, cần cấu hình Nginx chuyển tiếp địa chỉ và đặt \texttt{RAGTS\_\allowbreak{}TRUST\_\allowbreak{}PROXY\_\allowbreak{}HOPS} theo số proxy tin cậy thực tế.

Cơ chế này gắn với hệ thống giới hạn tần suất của backend. \texttt{ThrottlerGuard} khóa theo địa chỉ IP của client, và Express chỉ đọc được địa chỉ thật qua \texttt{req.\allowbreak{}ip} khi cấu hình \texttt{trust proxy} được đặt đúng. Trong \texttt{src/\allowbreak{}common/\allowbreak{}security.\allowbreak{}ts}:

\begin{minted}[fontsize=\footnotesize,breaklines,breakanywhere]{typescript}
export function trustProxySetting(nodeEnv: string): number | 'loopback' | false {
  if (nodeEnv === 'production') return 1; // one reverse proxy hop in front of the app
  if (nodeEnv === 'test') return 'loopback';
  return false; // dev: the client IP must stay the real one so `@Ip()` is meaningful
}
\end{minted}

Giá trị \texttt{1} ở production nghĩa là "một lớp proxy đáng tin cậy phía trước ứng dụng" \chapterSevenDash{} tức là Nginx. Hàm \texttt{app.\allowbreak{}setup.\allowbreak{}ts} đặt giá trị này \emph{trước} khi Express đọc địa chỉ client, với chú thích giải thích rõ: \texttt{trust proxy} phải được đặt trước vì rate limiter khóa theo địa chỉ đó.

Nếu đặt sai, có hai hậu quả nghiêm trọng. Nếu đặt quá nhiều lớp proxy tin cậy, kẻ tấn công có thể giả mạo địa chỉ qua header \texttt{X-\allowbreak{}Forwarded-\allowbreak{}For} và bỏ qua hoàn toàn giới hạn tần suất. Nếu đặt quá ít, tất cả người dùng đi qua cùng một proxy sẽ dùng chung một bộ đếm, khiến người dùng hợp pháp bị chặn vì hành vi của người khác.

Các mức giới hạn được định nghĩa tập trung trong \texttt{src/\allowbreak{}common/\allowbreak{}throttler.\allowbreak{}ts}:

\begingroup\small
\begin{longtable}{@{}>{\raggedright\arraybackslash}p{0.31\dimexpr\linewidth-4\tabcolsep\relax}>{\raggedright\arraybackslash}p{0.32\dimexpr\linewidth-4\tabcolsep\relax}>{\raggedright\arraybackslash}p{0.37\dimexpr\linewidth-4\tabcolsep\relax}@{}}
\toprule
\textbf{Tuyến} & \textbf{Giới hạn} & \textbf{Lý do} \\
\midrule
\endhead
\texttt{POST /\allowbreak{}auth/\allowbreak{}register} & 5 / 60 giây & Chống tạo tài khoản hàng loạt và làm cạn CPU bởi Argon2 \\
\texttt{POST /\allowbreak{}auth/\allowbreak{}login} & 10 / 60 giây & Chống đoán mật khẩu theo bộ từ điển \\
\texttt{POST /\allowbreak{}auth/\allowbreak{}refresh} & 30 / 60 giây & Hạn chế phát hành token hàng loạt \\
\texttt{POST /\allowbreak{}auth/\allowbreak{}change-\allowbreak{}password} & 5 / 60 giây & Hạn chế thay đổi mật khẩu liên tục \\
\texttt{POST /\allowbreak{}rag/\allowbreak{}chat}, \texttt{/\allowbreak{}rag/\allowbreak{}chat/\allowbreak{}stream} & 20 / 60 giây & LLM mặc định chỉ xử lý một lượt sinh tại một thời điểm \\
Mặc định (mọi tuyến khác) & 120 / 60 giây & Giới hạn mặc định cho các endpoint còn lại \\
\bottomrule
\end{longtable}
\endgroup

Chỉ một rate limiter tên \texttt{default} được đăng ký toàn cục, còn các tuyến nhạy cảm thu hẹp giới hạn của riêng mình bằng decorator \texttt{@\allowbreak{}RateLimit(.\allowbreak{}.\allowbreak{}.\allowbreak{})}. Đây là điểm thiết kế đáng chú ý: nếu đăng ký một rate limiter riêng cho mỗi tuyến, tất cả các rate limiter đó sẽ áp dụng lên mọi tuyến \chapterSevenDash{} điều không mong muốn. \texttt{ThrottlerGuard} được gắn qua \texttt{APP\_\allowbreak{}GUARD} nên kể cả các tuyến được thêm sau này cũng tự động được giới hạn, kể cả khi người đóng góp mới không nhớ gắn decorator.


\section{Cấu hình HTTPS và bảo mật kết nối}
\label{sec:cau-hinh-https}

\subsection{TLS termination tại Nginx}

Bản production của hệ thống thực hiện TLS termination ngay tại Nginx, theo đúng mô hình "kết thúc tại biên" (TLS termination at the edge): chứng thư được nạp một lần bởi Nginx, phần còn lại của hệ thống \chapterSevenDash{} từ Nginx đến container backend \chapterSevenDash{} chạy trong mạng nội bộ của Docker và không cần mã hóa.

Trong \texttt{nginx.\allowbreak{}prod.\allowbreak{}conf}, khối máy chủ ở cổng 443 khai báo \texttt{listen 443 ssl} cùng hai đường dẫn tới tệp chứng thư:

\begin{minted}[fontsize=\footnotesize,breaklines,breakanywhere]{nginx}
server {
    listen 443 ssl;
    server_name tenmien.com www.tenmien.com;

    ssl_certificate "/etc/nginx/ssl/cert.pem";
    ssl_certificate_key "/etc/nginx/ssl/key.pem";
    ...
}
\end{minted}

Hai đường dẫn này trỏ tới thư mục \texttt{/\allowbreak{}etc/\allowbreak{}nginx/\allowbreak{}ssl} bên trong container, là nơi thư mục \texttt{.\allowbreak{}/\allowbreak{}ssl} của máy chủ được gắn vào theo khai báo trong \texttt{docker-\allowbreak{}compose.\allowbreak{}prod.\allowbreak{}yaml}. Thư mục \texttt{ssl/\allowbreak{}} chứa hai tệp \texttt{cert.\allowbreak{}pem} và \texttt{key.\allowbreak{}pem} và đã được liệt kê trong \texttt{.\allowbreak{}gitignore} kể từ commit đầu tiên, cùng với \texttt{dist/\allowbreak{}}. Việc loại trừ này là bắt buộc: khóa riêng TLS không bao giờ được phép xuất hiện trong lịch sử kho mã nguồn, kể cả khi kho đó là kho nội bộ.

Khối máy chủ ở cổng 80 chuyển hướng sang địa chỉ HTTPS của cùng tên miền:

\begin{minted}[fontsize=\footnotesize,breaklines,breakanywhere]{nginx}
server {
    listen 80;
    server_name tenmien.com www.tenmien.com;
    return 301 https://$host$request_uri;
}
\end{minted}

Có một lý do kỹ thuật để giữ cả hai cổng: nhiều ứng dụng dựa trên web hiện đại, chẳng hạn trình duyệt nhúng, mong đợi chuyển hướng từ HTTP sang HTTPS. Chuyển hướng 301 giữ nguyên đường dẫn và chuỗi truy vấn, nên không phá vỡ các liên kết sâu vào trong ứng dụng.

Hai cấu hình cùng khai báo \texttt{server\_\allowbreak{}name tenmien.\allowbreak{}com www.\allowbreak{}tenmien.\allowbreak{}com} \chapterSevenDash{} bao gồm cả dạng không có tiền tố \texttt{www}. Cần lưu ý rằng giá trị \texttt{tenmien.\allowbreak{}com} trong tệp là giá trị thay thế, phải được thay bằng tên miền thật của đơn vị triển khai trước khi đưa vào phục vụ. Tương tự, đường dẫn chứng thư phải khớp với cặp tệp mà người vận hành đặt vào thư mục \texttt{ssl/\allowbreak{}}.

Một điểm cần lưu ý về mức độ mã hóa: vì khối \texttt{server} ở cổng 443 không khai báo tham số \texttt{ssl\_\allowbreak{}protocols} hay \texttt{ssl\_\allowbreak{}ciphers}, Nginx sẽ dùng bộ mặc định của bản \texttt{nginx:\allowbreak{}alpine} đang chạy. Điều này chấp nhận được ở mức cấu hình hiện tại, nhưng là hạn chế đã biết cần nêu trong Mục 7.8.10: nên chốt rõ \texttt{ssl\_\allowbreak{}protocols TLSv1.\allowbreak{}2 TLSv1.\allowbreak{}3} và \texttt{ssl\_\allowbreak{}prefer\_\allowbreak{}server\_\allowbreak{}ciphers off} khi đưa vào phục vụ thật.

\subsection{Các lớp bảo mật tích hợp trong backend}

Ngoài TLS ở tầng reverse proxy, backend còn áp dụng nhiều lớp bảo mật ở tầng ứng dụng. Việc hiểu rõ phạm vi của từng lớp giúp tránh nhầm lẫn khi xử lý sự cố.

\textbf{HTTP security header.} \texttt{helmet} được áp dụng cho mọi phản hồi trong \texttt{app.\allowbreak{}setup.\allowbreak{}ts}. Chính sách CSP được tùy biến theo đặc thù của API: \texttt{default-\allowbreak{}src 'self'}, \texttt{script-\allowbreak{}src 'self'}, \texttt{frameAncestors 'none'}, \texttt{object-\allowbreak{}src 'none'}. Cơ chế HSTS chỉ bật ở production với \texttt{maxAge:\allowbreak{} 31 536 000} (một năm) và \texttt{includeSubDomains:\allowbreak{} true}. Về chính sách cross-origin, \texttt{crossOriginEmbedderPolicy} bị tắt vì luồng SSE và công cụ GGUF cục bộ không cần cross-origin isolation.

\textbf{Ẩn tài liệu API ở production.} Backend chỉ bật Swagger khi \texttt{NODE\_\allowbreak{}ENV !=\allowbreak{}=\allowbreak{} 'production'}. Đây là biện pháp có chủ ý chứ không phải sự bỏ sót: tài liệu Swagger tiết lộ danh sách endpoint của hệ thống, bao gồm cả các endpoint quản trị, và là tài liệu tham khảo hữu ích cho kẻ thăm dò hệ thống.

\textbf{Giới hạn tần suất toàn cục.} Như đã trình bày ở Mục 7.6.3.

\textbf{Giới hạn thông tin trả về từ health endpoint.} \texttt{HealthController} có cơ chế che giấu thông tin riêng: ở production, phản hồi chỉ giữ lại trạng thái tổng quát, phiên bản, uptime và trạng thái của từng thành phần; chi tiết nhạy cảm như mã mô hình, thiết bị GPU và độ dài hàng đợi chỉ cung cấp cho quản trị viên qua \texttt{GET /\allowbreak{}api/\allowbreak{}v1/\allowbreak{}admin/\allowbreak{}health}.

\subsection{Quy trình kiểm tra trước khi đưa phiên bản mới vào phục vụ}

Trước khi đưa phiên bản mới vào phục vụ, người vận hành thực hiện theo thứ tự sau:

\begin{enumerate}
    \item Chạy \texttt{nginx -\allowbreak{}t} để kiểm tra cú pháp cấu hình.
    \item Nạp lại cấu hình.
    \item Dùng trình duyệt kiểm tra tên miền, chuỗi tin cậy cùng ngày hết hạn của chứng thư.
    \item Gửi câu hỏi qua chính địa chỉ HTTPS và quan sát luồng SSE đến khi hoàn tất.
    \item Kiểm tra phản hồi của \texttt{/\allowbreak{}api/\allowbreak{}v1/\allowbreak{}health/\allowbreak{}live} và \texttt{/\allowbreak{}api/\allowbreak{}v1/\allowbreak{}health/\allowbreak{}ready} qua chính kết nối HTTPS.
\end{enumerate}

Bước 4 và bước 5 phải thực hiện qua địa chỉ HTTPS chứ không phải qua địa chỉ cục bộ, vì nhiều cấu hình sai chỉ biểu hiện ở tầng TLS hoặc ở phía proxy.

Khi gia hạn chứng thư, cần nạp lại Nginx rồi đối chiếu ngày hết hạn mà máy khách thực sự nhận được \chapterSevenDash{} không chỉ đối chiếu với tệp trên đĩa, vì cấu hình có thể chưa được nạp lại.


\section{Giám sát, ghi log và quản lý hoạt động hệ thống}
\label{sec:giam-sat-logging}

\subsection{Các health endpoint của backend}

Hệ thống phân biệt ba mức kiểm tra, mỗi mức phục vụ một đối tượng khác nhau.

\textbf{Liveness (\texttt{/api/v1/health/live}).} Endpoint trả HTTP~200 khi tiến trình backend còn phản hồi. Docker dùng endpoint này cho \texttt{HEALTHCHECK} của container; load balancer cũng có thể dùng nó để kiểm tra kết nối tới backend. Endpoint không truy vấn MongoDB hay Qdrant. Nhờ đó, khi một dịch vụ phụ thuộc phản hồi chậm, kết quả liveness vẫn phản ánh đúng trạng thái của tiến trình backend.

\textbf{Readiness (\texttt{/\allowbreak{}api/\allowbreak{}v1/\allowbreak{}health/\allowbreak{}ready}).} Dựa trên kết quả thăm dò phụ thuộc. Nếu MongoDB hoặc Qdrant mất kết nối, tuyến trả HTTP 503. Nếu chỉ mô hình sinh chưa sẵn sàng, tuyến vẫn trả HTTP 200 với trạng thái \texttt{degraded}.

\textbf{Trạng thái tổng hợp (\texttt{/\allowbreak{}api/\allowbreak{}v1/\allowbreak{}health}).} Trả về đối tượng đầy đủ gồm trạng thái, phiên bản, uptime, phiên bản dữ liệu world và trạng thái của từng thành phần.

Quy tắc tổng hợp nằm trong hàm \texttt{overallStatus}:

\begin{minted}[fontsize=\footnotesize,breaklines,breakanywhere]{typescript}
export function overallStatus(c: Components): HealthStatus {
  if (!c.mongo || !c.qdrant) return 'down';
  if (!c.llm) return 'degraded';
  return 'ok';
}
\end{minted}

Phân biệt ba trạng thái này có ý nghĩa vận hành rõ ràng. \texttt{down} nghĩa là hệ thống không thể phục vụ chức năng cốt lõi vì mất kết nối cơ sở dữ liệu hoặc chỉ mục. \texttt{degraded} nghĩa là hệ thống vẫn phục vụ được nhưng chưa sinh câu trả lời được. \texttt{ok} nghĩa là mọi thành phần đều sẵn sàng.

Do vậy, để xác nhận chức năng hỏi đáp, cần đọc cả nội dung phản hồi chứ không chỉ mã HTTP, và phải thực hiện một lượt sinh thử. Một tuyến trả 200 với trạng thái \texttt{degraded} không đồng nghĩa với hệ thống sẵn sàng phục vụ.

\subsection{Cơ chế thăm dò nền và sự kiện thay đổi trạng thái}

\texttt{HealthService} không gọi phụ thuộc ngay khi nhận yêu cầu. Thay vào đó, nó chạy một vòng thăm dò nền mỗi 10 giây:

\begin{minted}[fontsize=\footnotesize,breaklines,breakanywhere]{typescript}
export const CHECK_INTERVAL_MS = 10_000;
\end{minted}

Mỗi yêu cầu HTTP chỉ đọc kết quả của lần thăm dò gần nhất. Cách tổ chức này có ba lợi ích rõ ràng: health endpoint không bao giờ trở thành điểm nghẽn khi tải cao; lượt thăm dò MongoDB được đặt thời hạn cứng 2 giây nên không treo; và các lời gọi thăm dò đồng thời được gộp lại bằng cơ chế \texttt{this.\allowbreak{}running}, tránh phát sinh nhiều lượt thăm dò chồng lên nhau.

Khi trạng thái một thành phần thay đổi, service ghi sự kiện \texttt{HEALTH\_\allowbreak{}CHANGED} kèm tên thành phần và chiều chuyển đổi:

\begin{minted}[fontsize=\footnotesize,breaklines,breakanywhere]{typescript}
this.logger.warn(
  { component: k, from: prev[k] ? 'up' : 'down', to: next[k] ? 'up' : 'down' },
  'HEALTH_CHANGED',
);
\end{minted}

Sự kiện này là công cụ chẩn đoán quan trọng khi người dùng báo "giao diện vẫn tải được nhưng hỏi đáp chậm". Người vận hành có thể đối chiếu sự kiện này với độ dài hàng đợi và thời điểm nhận token đầu tiên để xác định nguyên nhân nằm ở phụ thuộc hay ở hàng đợi suy luận.

Backend gọi \texttt{unref()} trên timer dùng cho health check. Khi các tác vụ khác đã hoàn tất, timer này không giữ Node.js event loop chạy thêm.

\subsection{Log có cấu trúc và cơ chế redaction}

Nhật ký được ghi bằng Pino qua \texttt{nestjs-\allowbreak{}pino}, với mỗi yêu cầu HTTP tương ứng đúng một dòng JSON.

Định dạng dòng nhật ký được chuẩn hoá theo hợp đồng API: \texttt{method}, \texttt{path} (không kèm chuỗi truy vấn), \texttt{status}, \texttt{user\_\allowbreak{}id}, cùng các trường ngữ cảnh do nghiệp vụ bổ sung. Log level được chọn theo mã trạng thái: \texttt{error} cho lỗi 5xx, \texttt{warn} cho lỗi 4xx, \texttt{info} cho các trường hợp còn lại. Riêng các sự kiện kết thúc bằng \texttt{\_\allowbreak{}ABORTED} được hạ xuống mức \texttt{info} vì một yêu cầu bị người dùng hủy không phải lỗi hệ thống.

Cơ chế che thông tin nhạy cảm được khai báo bằng danh sách \texttt{REDACT\_\allowbreak{}PATHS}:

\begin{minted}[fontsize=\footnotesize,breaklines,breakanywhere]{typescript}
export const REDACT_PATHS = [
  'req.headers.authorization',
  'req.headers.cookie',
  'password',
  'old_password',
  'new_password',
  'refresh_token',
  'access_token',
  'query',
  '*.password',
  '*.old_password',
  '*.new_password',
  '*.refresh_token',
  '*.access_token',
  '*.query',
];
\end{minted}

Giá trị bị thay bằng chuỗi \texttt{[redacted]}. Danh sách bao gồm cả header \texttt{Authorization} và \texttt{Cookie}, các biến thể mật khẩu, cả hai loại token, và trường \texttt{query} \chapterSevenDash{} tức là nội dung câu hỏi của người dùng. Việc không ghi nội dung truy vấn là quyết định bảo vệ quyền riêng tư: dữ liệu câu hỏi của người dùng không nên xuất hiện trong nhật ký hệ thống.

Đối với truy vấn của người dùng, log ghi độ dài qua trường \texttt{query\_len} để phục vụ chẩn đoán mà vẫn giữ kín nội dung truy vấn.

\subsection{Giám sát ở tầng Nginx}

Nhật ký ở tầng ứng dụng chỉ ghi những yêu cầu đã đi qua proxy và tới được backend. Những yêu cầu bị Nginx chặn trước \chapterSevenDash{} ví dụ bị từ chối vì giới hạn tần suất ở tầng proxy, hoặc gặp lỗi TLS \chapterSevenDash{} không xuất hiện trong nhật ký Pino. Vì vậy, giám sát đầy đủ đòi hỏi kết hợp hai nguồn.

Access log của Nginx, do cấu hình production ghi vào \texttt{/\allowbreak{}var/\allowbreak{}log/\allowbreak{}nginx/\allowbreak{}rag\_\allowbreak{}access.\allowbreak{}log}, cung cấp góc nhìn từ phía biên: địa chỉ client thật (nhờ \texttt{X-\allowbreak{}Real-\allowbreak{}IP} và \texttt{X-\allowbreak{}Forwarded-\allowbreak{}For}), mã trạng thái trả về, và thời gian xử lý. Đây là nguồn để đối chiếu khi người dùng báo "trang không tải được" \chapterSevenDash{} nếu không có dòng trong \texttt{rag\_\allowbreak{}access.\allowbreak{}log}, yêu cầu chưa từng tới Nginx, và vấn đề nằm ở tầng mạng hoặc DNS. Nếu có dòng với mã 502 hoặc 504, lỗi nằm ở đường Nginx đến backend.

Error log \texttt{rag\_\allowbreak{}error.\allowbreak{}log} ở mức \texttt{warn} ghi lại các sự kiện như backend không phản hồi trong \texttt{proxy\_\allowbreak{}read\_\allowbreak{}timeout} (dấu hiệu LLM quá tải), lỗi thư mục chứng thư, hoặc cấu hình không hợp lệ. Mức \texttt{warn} được chọn thay vì \texttt{error} nhằm lọc bỏ phần lớn thông báo thông thường của Nginx, khiến tệp này thực sự dùng để chẩn đoán.

Cả hai tệp nằm trong volume \texttt{prod-\allowbreak{}logs} được ánh xạ ra máy chủ, nên chúng tồn tại qua các lần khởi động lại container. Quy trình vận hành hằng ngày nên kiểm tra định kỳ ba nội dung: tỉ lệ yêu cầu có mã 5xx trong access log, số dòng trong error log, và việc tệp nhật ký có phình to bất thường do \texttt{try\_\allowbreak{}files} trả về 404 lặp lại.

\subsection{Request ID và cách truy vết lỗi}

Mỗi yêu cầu đều mang một mã định danh UUIDv7. Middleware \texttt{requestIdMiddleware} là middleware đầu tiên của ứng dụng: nó chấp nhận header \texttt{X-\allowbreak{}Request-\allowbreak{}Id} từ client chỉ khi giá trị đó là một UUID hợp lệ, nếu không thì tự sinh một UUIDv7 mới. Giá trị này được gán vào \texttt{req.\allowbreak{}id}, được gửi lại cho client trong response header \texttt{X-\allowbreak{}Request-\allowbreak{}Id}, và mở một ngữ cảnh yêu cầu thông qua \texttt{AsyncLocalStorage}.

Vì \texttt{AsyncLocalStorage} giữ ngữ cảnh theo luồng thực thi bất đồng bộ, mọi tầng nghiệp vụ có thể gắn trường vào dòng nhật ký hiện tại mà không cần truyền tham số. Hai hàm \texttt{logEvent(event, fields)} và \texttt{logFields(fields)} phục vụ mục đích này. Ngoài request context, phần thân lỗi của mọi endpoint cũng chứa trường \texttt{request\_\allowbreak{}id}, nên khi người dùng báo lỗi chỉ cần dán giá trị header \texttt{X-\allowbreak{}Request-\allowbreak{}Id} là tra được toàn bộ chuỗi thao tác tương ứng.

Việc chỉ chấp nhận UUID hợp lệ từ client có lý do bảo mật: nếu chấp nhận tùy ý, kẻ tấn công có thể chèn ký tự kiểm soát vào nhật ký và gây nhiễu khi phân tích.

\subsection{API quản trị và giám sát}

API quản trị được bảo vệ bởi ba lớp kiểm soát: \texttt{JwtAuthGuard} xác thực, \texttt{RolesGuard} với \texttt{@\allowbreak{}Roles('admin')} phân quyền, và kiểm tra lại ở tầng dịch vụ (\texttt{assertAdmin}) như một lớp phòng thủ sâu hơn.

\begingroup\small
\begin{longtable}{@{}>{\raggedright\arraybackslash}p{0.47\dimexpr\linewidth-2\tabcolsep\relax}>{\raggedright\arraybackslash}p{0.53\dimexpr\linewidth-2\tabcolsep\relax}@{}}
\toprule
\textbf{Tuyến} & \textbf{Mục đích} \\
\midrule
\endhead
\texttt{GET /\allowbreak{}admin/\allowbreak{}health} & Trạng thái chi tiết của từng thành phần, gồm mã mô hình và thiết bị \\
\texttt{GET /\allowbreak{}admin/\allowbreak{}stats} & Thống kê toàn hệ thống \\
\texttt{GET /\allowbreak{}admin/\allowbreak{}users} & Danh sách người dùng có phân trang \\
\texttt{GET /\allowbreak{}admin/\allowbreak{}users/\allowbreak{}:\allowbreak{}id} & Chi tiết một người dùng \\
\texttt{PATCH /\allowbreak{}admin/\allowbreak{}users/\allowbreak{}:\allowbreak{}id/\allowbreak{}roles} & Thay đổi vai trò \\
\texttt{POST /\allowbreak{}admin/\allowbreak{}users/\allowbreak{}:\allowbreak{}id/\allowbreak{}disable} & Khóa tài khoản \\
\texttt{POST /\allowbreak{}admin/\allowbreak{}users/\allowbreak{}:\allowbreak{}id/\allowbreak{}enable} & Mở khóa tài khoản \\
\texttt{POST /\allowbreak{}admin/\allowbreak{}users/\allowbreak{}:\allowbreak{}id/\allowbreak{}revoke-\allowbreak{}tokens} & Thu hồi toàn bộ token của người dùng \\
\texttt{GET /\allowbreak{}admin/\allowbreak{}users/\allowbreak{}:\allowbreak{}id/\allowbreak{}tokens} & Danh sách refresh token family đang hoạt động \\
\texttt{GET /\allowbreak{}admin/\allowbreak{}users/\allowbreak{}:\allowbreak{}id/\allowbreak{}sessions} & Phiên hội thoại của người dùng \\
\texttt{GET /\allowbreak{}admin/\allowbreak{}sessions/\allowbreak{}:\allowbreak{}sessionId} & Chi tiết một phiên kèm tin nhắn \\
\texttt{GET /\allowbreak{}admin/\allowbreak{}login-\allowbreak{}attempts} & Nhật ký đăng nhập thất bại gần đây \\
\texttt{GET /\allowbreak{}admin/\allowbreak{}audit-\allowbreak{}logs} & Nhật ký kiểm toán có phân trang \\
\bottomrule
\end{longtable}
\endgroup

API quản trị chủ yếu phục vụ thao tác đọc (read-mostly): không có đường dẫn nào xoá nội dung hội thoại hay xoá bản ghi kiểm toán. Bảng \texttt{admin\_\allowbreak{}audit\_\allowbreak{}logs} chỉ cho phép ghi thêm bản ghi mới (append-only) \chapterSevenDash{} không có đường dẫn cập nhật hay xoá, kể cả ở tầng dịch vụ. Điều này có nghĩa là mọi hành động quản trị, kể cả hành động chỉ đọc, đều để lại dấu vết với tác nhân, đối tượng, trạng thái trước và sau, địa chỉ IP và mã yêu cầu.

Có hai quy tắc bảo vệ chống tự khóa: quản trị viên không thể tự hạ quyền hoặc tự khóa tài khoản của chính mình (\texttt{409 SELF\_\allowbreak{}LOCKOUT\_\allowbreak{}FORBIDDEN}), và quản trị viên đang hoạt động cuối cùng không thể bị hạ quyền hoặc khóa (\texttt{409 LAST\_\allowbreak{}ADMIN\_\allowbreak{}PROTECTED}). Cả hai kiểm tra đều nằm ở tầng dịch vụ chứ không chỉ ở controller.

Điểm vận hành quan trọng: khi vai trò hoặc trạng thái tài khoản thay đổi, hệ thống phải gọi \texttt{AuthService.\allowbreak{}revokeAllForUser(userId)}, vừa thu hồi toàn bộ refresh token, vừa tăng \texttt{users.\allowbreak{}token\_\allowbreak{}version}. Claim \texttt{ver} trong access token mang giá trị \texttt{token\_\allowbreak{}version} tại thời điểm cấp; \texttt{JwtAuthGuard} từ chối token có \texttt{ver} không còn khớp. Nhờ vậy, access token đã cấp cho client mất hiệu lực ngay khi phiên bản token thay đổi, thay vì chờ hết hạn.

\subsection{Quy trình nâng cấp và rollback}

Khi nâng cấp, cần tuân thủ trình tự sau:

\begin{enumerate}
    \item Ghi lại phiên bản image ứng dụng, tên mô hình và phiên bản dữ liệu tương ứng.
    \item Sao lưu MongoDB trước các thay đổi cấu trúc.
    \item Thay image ứng dụng.
    \item Lặp lại phép thử hỏi đáp có dẫn nguồn tại Mục 7.4.3.
\end{enumerate}

Thông tin phiên bản ở bước 1 là cơ sở để rollback. Khi cần khôi phục, dùng lại bộ phiên bản tương thích và kiểm tra cả khả năng đọc lịch sử hội thoại lẫn truy hồi bài báo. Bước kiểm tra lịch sử hội thoại là bắt buộc vì lịch sử được lưu bằng khoá UUIDv7 và có cột \texttt{deleted\_\allowbreak{}at} cho xoá mềm; một thay đổi cấu trúc lược đồ không đúng có thể làm mất khả năng đọc dữ liệu cũ mà không gây lỗi ngay lập tức.

\subsection{Quy trình chứng nhận hiệu năng sau thay đổi}

Mọi thay đổi lớn ở tầng nghiệp vụ cần được kiểm chứng bằng bộ kiểm thử tải ba tầng trong thư mục \texttt{benchmark/\allowbreak{}}, theo Quy tắc L1: một mức tải chỉ được công nhận chính thức nếu hệ thống duy trì được mức tải đó liên tục trong 180 giây mà giữ vững \texttt{P95 < 500 ms}, \texttt{P99 < 1000 ms} và tỉ lệ lỗi dưới 1\%.

\begingroup\small
\begin{longtable}{@{}>{\raggedright\arraybackslash}p{0.31\dimexpr\linewidth-4\tabcolsep\relax}>{\raggedright\arraybackslash}p{0.32\dimexpr\linewidth-4\tabcolsep\relax}>{\raggedright\arraybackslash}p{0.37\dimexpr\linewidth-4\tabcolsep\relax}@{}}
\toprule
\textbf{Tầng} & \textbf{Kịch bản} & \textbf{Mục tiêu} \\
\midrule
\endhead
A \chapterSevenDash{} Xác định ngưỡng quá tải & \texttt{k6\_\allowbreak{}prod\_\allowbreak{}open\_\allowbreak{}breakpoint.\allowbreak{}ts} & Mô hình tải mở (\texttt{ramping-\allowbreak{}arrival-\allowbreak{}rate}), quét từ 50 đến 500 RPS trên 6 tuyến đọc/ghi, xác định mức tải bắt đầu vượt ngưỡng cho phép \\
B \chapterSevenDash{} Kiểm chứng khả năng duy trì tải & \texttt{k6\_\allowbreak{}prod\_\allowbreak{}certified\_\allowbreak{}holdout.\allowbreak{}ts} + \texttt{collect\_\allowbreak{}telemetry.\allowbreak{}js} & Giữ tải 3 phút tại mức mục tiêu, thu thập chỉ số CPU/RAM mỗi 2 giây \\
C \chapterSevenDash{} Đo riêng chi phí xử lý mật mã & \texttt{k6\_\allowbreak{}prod\_\allowbreak{}auth\_\allowbreak{}crypto.\allowbreak{}ts} & Tải đồng thời \texttt{POST /\allowbreak{}auth/\allowbreak{}login} (Argon2id 64 MB), \texttt{POST /\allowbreak{}auth/\allowbreak{}refresh} và dò tình trạng kẹt event loop \\
\bottomrule
\end{longtable}
\endgroup

Mô hình tải mở (\texttt{ramping-\allowbreak{}arrival-\allowbreak{}rate}) được dùng thay cho mô hình tải đóng có sẵn vì lý do kỹ thuật đã được ghi nhận: trong mô hình tải đóng, khi máy chủ chậm đi, các người dùng ảo cũng tự động gửi ít yêu cầu hơn, làm đẹp số liệu giả tạo \chapterSevenDash{} hiện tượng được gọi là coordinated omission.

Kết quả chứng nhận hiện tại của hệ thống, đo trong 3 phút duy trì tải liên tục theo Quy tắc L1:

\begingroup\small
\begin{longtable}{@{}>{\raggedright\arraybackslash}p{0.22\dimexpr\linewidth-6\tabcolsep\relax}>{\raggedleft\arraybackslash}p{0.18\dimexpr\linewidth-6\tabcolsep\relax}>{\raggedleft\arraybackslash}p{0.31\dimexpr\linewidth-6\tabcolsep\relax}>{\raggedleft\arraybackslash}p{0.29\dimexpr\linewidth-6\tabcolsep\relax}@{}}
\toprule
\textbf{Chỉ số} & \textbf{Trước tối ưu} & \textbf{Sau tối ưu} & \textbf{Cải thiện} \\
\midrule
\endhead
Số người dùng đồng thời & 100 CUs (65 RPS) & 400 CUs (250 RPS) & 4,0 lần \\
Tổng yêu cầu trong 3 phút & 12.869 & 48.874 & 3,8 lần \\
Độ trễ trung vị (P50) & 160,29 ms & 29,79 ms & nhanh hơn 5,38 lần \\
Độ trễ phân vị 95 (P95) & 259,68 ms & 265,54 ms & đạt ngưỡng < 500 ms \\
Độ trễ phân vị 99 (P99) & 353,35 ms & 475,11 ms & đạt ngưỡng < 1000 ms \\
Tỉ lệ lỗi HTTP & 0,00\% & 0,00\% & 100\% \\
\bottomrule
\end{longtable}
\endgroup

Điểm đáng chú ý trong bảng này là độ trễ P95 gần như không đổi (259,68 ms so với 265,54 ms) trong khi mức tải tăng gấp bốn lần. Đây là kết quả của việc tối ưu ở tầng kiến trúc \chapterSevenDash{} bộ nhớ đệm \texttt{LRUCache} có gắn thế hệ theo tài khoản, cơ chế \texttt{Singleflight} gộp các yêu cầu thăm dò trùng lặp, truy vấn độc lập chạy song song, và \texttt{.\allowbreak{}lean()} bỏ qua chi phí hydrate của Mongoose \chapterSevenDash{} chứ không phải kết quả của việc tăng tài nguyên máy chủ.

Bốn yếu tố trong số liệu này có ý nghĩa vận hành riêng biệt và nên được ghi nhớ:

\texttt{UV\_\allowbreak{}THREADPOOL\_\allowbreak{}SIZE =\allowbreak{} 4} được đặt ngay dòng đầu tiên của \texttt{src/\allowbreak{}main.\allowbreak{}ts}, trước cả khi nạp \texttt{dotenv}. Con số này phải khớp với \texttt{parallelism:\allowbreak{} 4} của Argon2id trong cấu hình băm mật khẩu. Thử nghiệm tăng lên 12 đã gây tranh chấp băng thông DRAM và làm độ trễ tăng 26,7\%. Dòng đặt này phải đứng trước mọi lệnh khác vì thư viện native đọc giá trị này một lần lúc nạp.

Cơ chế ngăn cache stampede khi cache còn trống: khi hơn 100 người dùng ảo đồng thời yêu cầu \texttt{GET /\allowbreak{}chat-\allowbreak{}sessions/\allowbreak{}:\allowbreak{}id} khi cache chưa có dữ liệu, nhiều yêu cầu sẽ gặp cache miss trong cùng một mili giây và phát sinh các cặp truy vấn MongoDB trùng lặp qua đường hầm. Cơ chế \texttt{Singleflight} gộp các lời gọi trùng lặp thành một lần truy vấn duy nhất.

Hàng đợi suy luận được khảo sát riêng bằng \texttt{k6\_\allowbreak{}queue\_\allowbreak{}experiment.\allowbreak{}ts} để đánh giá cơ chế xếp hàng và hành vi trả mã 429 khi bão hòa.

Bộ kiểm thử tích hợp với dữ liệu thật nạp sẵn 1.500 phiên và 15.000 tin nhắn có trích dẫn, đoạn truy hồi và thời gian xử lý, tạo ra phản hồi khoảng 4,6 KB mỗi lần \chapterSevenDash{} gần với điều kiện thực tế thay vì truy vấn bảng rỗng.

\subsection{Xử lý sự cố thường gặp}

Bảng dưới đây tổng hợp các triệu chứng đã được ghi nhận trong tài liệu vận hành, kèm nguyên nhân và cách xử lý.

\begingroup\small
\begin{longtable}{@{}>{\raggedright\arraybackslash}p{0.47\dimexpr\linewidth-2\tabcolsep\relax}>{\raggedright\arraybackslash}p{0.53\dimexpr\linewidth-2\tabcolsep\relax}@{}}
\toprule
\textbf{Hiện tượng} & \textbf{Nguyên nhân / Cách xử lý} \\
\midrule
\endhead
Backend dừng ngay, in \texttt{Invalid environment} & Thiếu biến trong tệp môi trường; thông điệp liệt kê đủ tên biến sai \\
\texttt{components.\allowbreak{}mongo:\allowbreak{} false}, nhật ký \texttt{Authentication failed} & Sai tên đăng nhập hoặc mật khẩu, hoặc thiếu \texttt{authSource=\allowbreak{}admin} trong URI \\
\texttt{components.\allowbreak{}mongo:\allowbreak{} false}, không kết nối được & Tiến trình \texttt{cloudflared access tcp} chưa chạy hoặc đã tắt \\
\texttt{components.\allowbreak{}qdrant:\allowbreak{} false}, hỏi đáp trả 503 \texttt{qdrant unavailable} & Sai \texttt{QDRANT\_\allowbreak{}URL} hoặc \texttt{QDRANT\_\allowbreak{}API\_\allowbreak{}KEY}, hoặc máy chủ đang tắt \\
\texttt{components.\allowbreak{}llm:\allowbreak{} false}, hỏi đáp trả 503 \texttt{llm unavailable} & Sai \texttt{MODEL\_\allowbreak{}PATH}, hoặc mô hình chưa nạp xong (cần chờ vài giây) \\
Cảnh báo \texttt{QDRANT\_\allowbreak{}DIMENSION\_\allowbreak{}MISMATCH} & Mô hình embedding và chỉ mục dùng số chiều khác nhau; cần lập chỉ mục lại \\
Cảnh báo \texttt{QDRANT\_\allowbreak{}NO\_\allowbreak{}TWIG\_\allowbreak{}ID} & Chỉ mục chưa có trường \texttt{twig\_\allowbreak{}id}; chức năng so sánh không giới hạn theo nhánh \\
Hỏi đáp trả 429 \texttt{llm queue is full} & Hàng đợi LLM đầy; giảm tải hoặc tăng \texttt{LLM\_\allowbreak{}MAX\_\allowbreak{}QUEUE} \\
Hỏi đáp trả 504 \texttt{llm timeout} & Lượt sinh vượt \texttt{LLM\_\allowbreak{}TIMEOUT\_\allowbreak{}MS} \\
Giao diện báo "Backend không phản hồi" & Backend chưa chạy ở cổng 3000, hoặc đang chạy \texttt{npm run dev} thay vì \texttt{npm run dev:\allowbreak{}live} \\
\texttt{Port 5174 is already in use} & Có tiến trình khác chiếm cổng; chạy \texttt{npm run dev:\allowbreak{}live -\allowbreak{}-\allowbreak{} -\allowbreak{}-\allowbreak{}port 5175} \\
Giao diện vẫn dùng dữ liệu giả lập sau khi đổi \texttt{.\allowbreak{}env} & Chưa build lại frontend; Vite ghi giá trị biến vào bundle khi build \\
Mọi yêu cầu API trả 502 qua Nginx & Tên service trong \texttt{nginx.\allowbreak{}*.\allowbreak{}conf} không khớp với tên service trong Compose, hoặc backend chưa sẵn sàng \\
Nginx không khởi động, báo lỗi cấu hình & Thiếu \texttt{ssl/\allowbreak{}cert.\allowbreak{}pem} hoặc \texttt{ssl/\allowbreak{}key.\allowbreak{}pem} trong thư mục \texttt{ssl/\allowbreak{}} mà \texttt{docker-\allowbreak{}compose.\allowbreak{}prod.\allowbreak{}yaml} ánh xạ \\
Giao diện trả về 404 sau khi triển khai lại & Thư mục \texttt{frontend-\allowbreak{}rag/\allowbreak{}dist} chưa được dựng lại, hoặc \texttt{try\_\allowbreak{}files} không trỏ tới \texttt{/\allowbreak{}index.\allowbreak{}html} \\
Luồng SSE nhận câu trả lời theo từng đoạn lớn thay vì theo dòng & Thiếu \texttt{proxy\_\allowbreak{}buffering off} trong khối \texttt{location /\allowbreak{}api/\allowbreak{}} \\
Mọi người dùng bị giới hạn tần suất cùng lúc & Thiếu header \texttt{X-\allowbreak{}Forwarded-\allowbreak{}For} hoặc \texttt{trust proxy} chưa được đặt đúng số lớp proxy \\
\bottomrule
\end{longtable}
\endgroup

Một nhận xét xuyên suốt bảng này: phần lớn các lỗi vận hành đều không phải lỗi hệ thống mà là lỗi cấu hình môi trường, và gần như tất cả đều để lại dấu vết trong nhật ký hệ thống. Đây là hệ quả trực tiếp của việc kiểm tra cấu hình tại thời điểm khởi động và của cơ chế ghi sự kiện có cấu trúc: người vận hành không phải đoán nguyên nhân, mà chỉ cần đọc dòng nhật ký tương ứng.

\subsection{Giới hạn đã biết và hướng cải tiến vận hành}

Cuối chương, cần ghi nhận những hạn chế còn tồn tại trong quy trình vận hành hiện tại, vì chúng quyết định hướng phát triển của Chương 8.

Thứ nhất, \textbf{giới hạn của hàng đợi suy luận khi ngắt kết nối}. Slot suy luận của LLM chỉ được giải phóng khi sinh tới token đầu tiên, tức khoảng 15 giây với \texttt{top\_\allowbreak{}k} là 5 trên nền Metal. Trong khoảng thời gian đó, một lượt đã hủy vẫn chiếm chỗ. Đây là hệ quả không tránh được của việc xử lý prompt dài để tạo KV cache cho mô hình.

Thứ hai, \textbf{khoảng cách giữa môi trường phát triển và môi trường thật}. Docker trên macOS không có Metal nên LLM chạy trên CPU; mô hình 7B cần khoảng 6 GB RAM cho container. Máy ảo Docker cấp phát 8 GB kèm container khác có thể bị dừng với exit code 137.

Thứ ba, \textbf{sự phụ thuộc vào dữ liệu của pipeline bên ngoài}. Khi pipeline chưa xuất tệp world và các trường payload \texttt{twig\_\allowbreak{}id}, \texttt{chunk\_\allowbreak{}index}, sự kiện \texttt{route} trả về rỗng và chức năng so sánh không lọc theo nhánh. Đây là hạn chế về mặt dữ liệu chứ không phải về mặt triển khai.

Thứ tư, \textbf{chưa có giám sát tập trung}. Hiện tại nhật ký được ghi ra stdout theo định dạng JSON một dòng, phù hợp để thu thập bằng \texttt{docker logs} hoặc đẩy vào một hệ thống tổng hợp log, nhưng dự án chưa cung cấp cấu hình sẵn cho một hệ thống giám sát tập trung. Cùng với đó, việc quản lý chứng thư TLS chưa được tự động hoá.

Thứ năm, \textbf{cấu hình Nginx ở tầng biên chưa được đồng bộ với tầng dữ liệu}. Như đã phân tích ở Mục 7.2.4, tầng biên dùng tên service \texttt{backend-\allowbreak{}rag} và cổng 8000 \chapterSevenDash{} bộ thông số thuộc về giai đoạn backend dùng LanceDB \chapterSevenDash{} trong khi tầng dữ liệu hiện tại đã chuyển sang Qdrant với tên service \texttt{app} và cổng 3000. Nếu chạy hai tầng Compose cùng lúc, mọi yêu cầu API qua Nginx sẽ trả 502 do không phân giải được tên service. Đây là hạn chế cần khắc phục trước khi triển khai production thực tế, và cách khắc phục tự nhiên nhất là đồng bộ tên service cùng cổng giữa hai tầng, kèm cập nhật khối \texttt{upstream} trong cả \texttt{nginx.\allowbreak{}dev.\allowbreak{}conf} lẫn \texttt{nginx.\allowbreak{}prod.\allowbreak{}conf}.

Thứ sáu, \textbf{chưa chốt tham số TLS và chưa có tự động luân chuyển chứng thư}. Khối máy chủ ở cổng 443 không khai báo \texttt{ssl\_\allowbreak{}protocols} hay \texttt{ssl\_\allowbreak{}ciphers}, nên Nginx dùng cấu hình mặc định của image \texttt{nginx:\allowbreak{}alpine}. Bộ mặc định này vẫn chấp nhận được, nhưng để có hành vi xác định và không lệ thuộc vào phiên bản image, nên chốt rõ \texttt{ssl\_\allowbreak{}protocols TLSv1.\allowbreak{}2 TLSv1.\allowbreak{}3}. Bên cạnh đó, việc gia hạn chứng thư vẫn là thao tác thủ công: tệp trong thư mục \texttt{ssl/\allowbreak{}} được thay thế rồi container Nginx nạp lại, và không có cơ chế tự động nào theo dõi thời hạn.

Thứ bảy, \textbf{khối \texttt{upstream} mới chỉ có một bản sao}. Tham số \texttt{max\_\allowbreak{}fails} và \texttt{fail\_\allowbreak{}timeout} đã được cấu hình trong \texttt{nginx.\allowbreak{}prod.\allowbreak{}conf}, nhưng chỉ có một dòng \texttt{server}. Do đó, cơ chế \texttt{proxy\_\allowbreak{}next\_\allowbreak{}upstream} \chapterSevenDash{} vốn được khai báo đầy đủ danh sách mã lỗi \chapterSevenDash{} hiện chưa có tác dụng thực tế. Muốn có khả năng chịu lỗi thật sự, cần bổ sung các bản sao backend, và khi đó cân nhắc thêm chia tải theo thuật toán băm để phân bổ đều tải.

Thứ tám, \textbf{việc quản lý secret còn thủ công}. Hai khóa ký JWT và khóa truy cập Qdrant được cấp qua biến môi trường của Compose hoặc tệp \texttt{.\allowbreak{}env} cục bộ. Cơ chế chặn giá trị placeholder trong production đã ngăn được lỗi phổ biến nhất, nhưng việc luân chuyển khóa, thu hồi và lưu trữ trong hệ thống quản lý secret vẫn là thao tác thủ công. Cùng cơ chế đó, khóa riêng TLS trong thư mục \texttt{ssl/\allowbreak{}} cũng được quản lý theo cách tương tự \chapterSevenDash{} tệp nằm ngoài kho mã nguồn, nhưng không có hệ thống quản lý secret nào quản lý vòng đời của nó.


\section*{Tóm tắt Chương 7}
\addcontentsline{toc}{section}{Tóm tắt Chương 7}

Chương này đã trình bày chu trình vận hành đầy đủ của RAG\_TS, từ chuẩn bị môi trường đến giám sát sau khi triển khai. Năm kết luận chính có thể rút ra:

\textbf{Thứ nhất, kiểm tra cấu hình tại thời điểm khởi động là cơ chế vận hành quan trọng nhất.} Lược đồ Zod với các quy tắc siết chặt ở production, cộng với cú pháp \texttt{\$\{VAR:\allowbreak{}?message\}} của Compose, tạo thành hai lớp phòng thủ độc lập trước những cấu hình sai. Điều này đặc biệt có giá trị với hệ thống tự triển khai, nơi phần lớn lỗi vận hành là lỗi cấu hình chứ không phải lỗi logic.

\textbf{Thứ hai, đóng gói hai giai đoạn với bind mount chỉ đọc là biện pháp bảo đảm tính đúng đắn, không chỉ là biện pháp an toàn.} Cờ \texttt{:\allowbreak{}ro} trên \texttt{/\allowbreak{}models} và \texttt{/\allowbreak{}data} cùng với quy ước chỉ đọc đối với chỉ mục vector bảo đảm rằng tính toàn vẹn giữa chỉ mục, cây tri thức và kết quả khai phá được giữ xuyên suốt vòng đời triển khai.

\textbf{Thứ ba, mã lỗi của API giúp khoanh vùng sự cố.} Sự phân biệt giữa \texttt{503 llm unavailable} (chưa sẵn sàng), \texttt{429 llm queue is full} (quá tải), \texttt{504 llm timeout} (vượt thời gian) và \texttt{503 qdrant unavailable} (mất kết nối tới Qdrant) cho phép người vận hành xác định nguyên nhân từ một dòng nhật ký, thay vì phải dò theo từng thành phần.

\textbf{Thứ tư, mọi thay đổi triển khai đều phải kèm kiểm chứng hiệu năng theo Quy tắc L1.} Bộ kiểm thử ba tầng với mô hình mở để tránh coordinated omission đã công nhận mức tải 400 người dùng đồng thời với độ trễ P95 là 265,54 ms và tỉ lệ lỗi 0,00\% trong suốt 180 giây \chapterSevenDash{} con số đạt được nhờ tối ưu kiến trúc chứ không phải nhờ bổ sung tài nguyên.

\textbf{Thứ năm, Nginx có cấu hình và vòng đời triển khai riêng.} Hai tệp \texttt{nginx.dev.conf}, \texttt{nginx.prod.conf} và hai tệp Compose ở thư mục gốc quy định cách phục vụ tệp tĩnh, xử lý client-side routing, chuyển tiếp HTTP header và thực hiện TLS termination. Cấu hình này ảnh hưởng trực tiếp đến hành vi của backend: \texttt{proxy\_buffering off} giúp chuyển từng sự kiện SSE tới trình duyệt, còn \texttt{X-Forwarded-For} cung cấp địa chỉ client để áp dụng rate limit.
\par\endgroup
